% Options for packages loaded elsewhere
% Options for packages loaded elsewhere
\PassOptionsToPackage{unicode}{hyperref}
\PassOptionsToPackage{hyphens}{url}
\PassOptionsToPackage{dvipsnames,svgnames,x11names}{xcolor}
%
\documentclass[
  11pt,
  letterpaper,
  DIV=11,
  numbers=noendperiod]{scrreprt}
\usepackage{xcolor}
\usepackage[top=1in,bottom=1in,left=1in,right=1in]{geometry}
\usepackage{amsmath,amssymb}
\setcounter{secnumdepth}{3}
\usepackage{iftex}
\ifPDFTeX
  \usepackage[T1]{fontenc}
  \usepackage[utf8]{inputenc}
  \usepackage{textcomp} % provide euro and other symbols
\else % if luatex or xetex
  \usepackage{unicode-math} % this also loads fontspec
  \defaultfontfeatures{Scale=MatchLowercase}
  \defaultfontfeatures[\rmfamily]{Ligatures=TeX,Scale=1}
\fi
\usepackage{lmodern}
\ifPDFTeX\else
  % xetex/luatex font selection
\fi
% Use upquote if available, for straight quotes in verbatim environments
\IfFileExists{upquote.sty}{\usepackage{upquote}}{}
\IfFileExists{microtype.sty}{% use microtype if available
  \usepackage[]{microtype}
  \UseMicrotypeSet[protrusion]{basicmath} % disable protrusion for tt fonts
}{}
\makeatletter
\@ifundefined{KOMAClassName}{% if non-KOMA class
  \IfFileExists{parskip.sty}{%
    \usepackage{parskip}
  }{% else
    \setlength{\parindent}{0pt}
    \setlength{\parskip}{6pt plus 2pt minus 1pt}}
}{% if KOMA class
  \KOMAoptions{parskip=half}}
\makeatother
% Make \paragraph and \subparagraph free-standing
\makeatletter
\ifx\paragraph\undefined\else
  \let\oldparagraph\paragraph
  \renewcommand{\paragraph}{
    \@ifstar
      \xxxParagraphStar
      \xxxParagraphNoStar
  }
  \newcommand{\xxxParagraphStar}[1]{\oldparagraph*{#1}\mbox{}}
  \newcommand{\xxxParagraphNoStar}[1]{\oldparagraph{#1}\mbox{}}
\fi
\ifx\subparagraph\undefined\else
  \let\oldsubparagraph\subparagraph
  \renewcommand{\subparagraph}{
    \@ifstar
      \xxxSubParagraphStar
      \xxxSubParagraphNoStar
  }
  \newcommand{\xxxSubParagraphStar}[1]{\oldsubparagraph*{#1}\mbox{}}
  \newcommand{\xxxSubParagraphNoStar}[1]{\oldsubparagraph{#1}\mbox{}}
\fi
\makeatother


\usepackage{longtable,booktabs,array}
\usepackage{calc} % for calculating minipage widths
% Correct order of tables after \paragraph or \subparagraph
\usepackage{etoolbox}
\makeatletter
\patchcmd\longtable{\par}{\if@noskipsec\mbox{}\fi\par}{}{}
\makeatother
% Allow footnotes in longtable head/foot
\IfFileExists{footnotehyper.sty}{\usepackage{footnotehyper}}{\usepackage{footnote}}
\makesavenoteenv{longtable}
\usepackage{graphicx}
\makeatletter
\newsavebox\pandoc@box
\newcommand*\pandocbounded[1]{% scales image to fit in text height/width
  \sbox\pandoc@box{#1}%
  \Gscale@div\@tempa{\textheight}{\dimexpr\ht\pandoc@box+\dp\pandoc@box\relax}%
  \Gscale@div\@tempb{\linewidth}{\wd\pandoc@box}%
  \ifdim\@tempb\p@<\@tempa\p@\let\@tempa\@tempb\fi% select the smaller of both
  \ifdim\@tempa\p@<\p@\scalebox{\@tempa}{\usebox\pandoc@box}%
  \else\usebox{\pandoc@box}%
  \fi%
}
% Set default figure placement to htbp
\def\fps@figure{htbp}
\makeatother


% definitions for citeproc citations
\NewDocumentCommand\citeproctext{}{}
\NewDocumentCommand\citeproc{mm}{%
  \begingroup\def\citeproctext{#2}\cite{#1}\endgroup}
\makeatletter
 % allow citations to break across lines
 \let\@cite@ofmt\@firstofone
 % avoid brackets around text for \cite:
 \def\@biblabel#1{}
 \def\@cite#1#2{{#1\if@tempswa , #2\fi}}
\makeatother
\newlength{\cslhangindent}
\setlength{\cslhangindent}{1.5em}
\newlength{\csllabelwidth}
\setlength{\csllabelwidth}{3em}
\newenvironment{CSLReferences}[2] % #1 hanging-indent, #2 entry-spacing
 {\begin{list}{}{%
  \setlength{\itemindent}{0pt}
  \setlength{\leftmargin}{0pt}
  \setlength{\parsep}{0pt}
  % turn on hanging indent if param 1 is 1
  \ifodd #1
   \setlength{\leftmargin}{\cslhangindent}
   \setlength{\itemindent}{-1\cslhangindent}
  \fi
  % set entry spacing
  \setlength{\itemsep}{#2\baselineskip}}}
 {\end{list}}
\usepackage{calc}
\newcommand{\CSLBlock}[1]{\hfill\break\parbox[t]{\linewidth}{\strut\ignorespaces#1\strut}}
\newcommand{\CSLLeftMargin}[1]{\parbox[t]{\csllabelwidth}{\strut#1\strut}}
\newcommand{\CSLRightInline}[1]{\parbox[t]{\linewidth - \csllabelwidth}{\strut#1\strut}}
\newcommand{\CSLIndent}[1]{\hspace{\cslhangindent}#1}



\setlength{\emergencystretch}{3em} % prevent overfull lines

\providecommand{\tightlist}{%
  \setlength{\itemsep}{0pt}\setlength{\parskip}{0pt}}



 


% --------------------------------------------------
% FastMAPOL ATBD PDF formatting
% --------------------------------------------------

\usepackage{scrlayer-scrpage}
\usepackage{xcolor}
\usepackage{booktabs}
\usepackage{longtable}
\usepackage{array}
\usepackage{microtype}

% --------------------------------------------------
% Colors
% --------------------------------------------------

\definecolor{NASABlue}{HTML}{0B3D91}
\definecolor{NASARed}{HTML}{FC3D21}
\definecolor{DarkGray}{HTML}{444444}

% --------------------------------------------------
% Headings
% --------------------------------------------------

\addtokomafont{chapter}{\color{NASABlue}}
\addtokomafont{section}{\color{NASABlue}}
\addtokomafont{subsection}{\color{DarkGray}}

% Slightly cleaner chapter heading
\RedeclareSectionCommand[
  beforeskip=0pt,
  afterskip=1.5\baselineskip
]{chapter}

% --------------------------------------------------
% Header / footer
% --------------------------------------------------

\clearpairofpagestyles

\ihead{\small FastMAPOL ATBD}
\ohead{\small PACE Mission}

\ifoot{\small NASA GSFC}
\ofoot{\small \pagemark}

\setkomafont{pageheadfoot}{\normalfont\small\color{DarkGray}}
\setkomafont{pagenumber}{\normalfont\small}

\pagestyle{scrheadings}

% --------------------------------------------------
% Tables
% --------------------------------------------------

\renewcommand{\arraystretch}{1.15}

% --------------------------------------------------
% Paragraph spacing
% --------------------------------------------------

\setlength{\parindent}{0pt}
\setlength{\parskip}{0.45em}
\KOMAoption{captions}{tableheading}
\makeatletter
\@ifpackageloaded{bookmark}{}{\usepackage{bookmark}}
\makeatother
\makeatletter
\@ifpackageloaded{caption}{}{\usepackage{caption}}
\AtBeginDocument{%
\ifdefined\contentsname
  \renewcommand*\contentsname{Table of contents}
\else
  \newcommand\contentsname{Table of contents}
\fi
\ifdefined\listfigurename
  \renewcommand*\listfigurename{List of Figures}
\else
  \newcommand\listfigurename{List of Figures}
\fi
\ifdefined\listtablename
  \renewcommand*\listtablename{List of Tables}
\else
  \newcommand\listtablename{List of Tables}
\fi
\ifdefined\figurename
  \renewcommand*\figurename{Figure}
\else
  \newcommand\figurename{Figure}
\fi
\ifdefined\tablename
  \renewcommand*\tablename{Table}
\else
  \newcommand\tablename{Table}
\fi
}
\@ifpackageloaded{float}{}{\usepackage{float}}
\floatstyle{ruled}
\@ifundefined{c@chapter}{\newfloat{codelisting}{h}{lop}}{\newfloat{codelisting}{h}{lop}[chapter]}
\floatname{codelisting}{Listing}
\newcommand*\listoflistings{\listof{codelisting}{List of Listings}}
\makeatother
\makeatletter
\makeatother
\makeatletter
\@ifpackageloaded{caption}{}{\usepackage{caption}}
\@ifpackageloaded{subcaption}{}{\usepackage{subcaption}}
\makeatother
\usepackage{bookmark}
\IfFileExists{xurl.sty}{\usepackage{xurl}}{} % add URL line breaks if available
\urlstyle{same}
\hypersetup{
  pdftitle={NASA FastMAPOL Algorithm Theoretical Basis Document},
  pdfauthor={Meng Gao, Pengwang Zhai, Kamal Aryal, Kirk Knobelspiesse, Bryan Franz},
  colorlinks=true,
  linkcolor={blue},
  filecolor={Maroon},
  citecolor={blue},
  urlcolor={blue},
  pdfcreator={LaTeX via pandoc}}


\title{NASA FastMAPOL Algorithm Theoretical Basis Document}
\usepackage{etoolbox}
\makeatletter
\providecommand{\subtitle}[1]{% add subtitle to \maketitle
  \apptocmd{\@title}{\par {\large #1 \par}}{}{}
}
\makeatother
\subtitle{PACE Polarimeter Aerosol and Surface Retrieval}
\author{Meng Gao, Pengwang Zhai, Kamal Aryal, Kirk Knobelspiesse, Bryan
Franz}
\date{2026-08-26}
\begin{document}
\maketitle

\renewcommand*\contentsname{Table of contents}
{
\hypersetup{linkcolor=}
\setcounter{tocdepth}{2}
\tableofcontents
}
\listoffigures
\listoftables

\bookmarksetup{startatroot}

\chapter{FastMAPOL Algorithm Theoretical Basis
Document}\label{fastmapol-algorithm-theoretical-basis-document}

PACE Multi-Angle Polarimetric Aerosol, Ocean, and Land Retrieval

\hfill\break

\bookmarksetup{startatroot}

\chapter{Executive Summary}\label{executive-summary}

FastMAPOL (Fast Multi-Angle Polarimetric Ocean and Land) is a nonlinear
least-squares retrieval algorithm developed for the NASA Plankton,
Aerosol, Cloud, ocean Ecosystem (PACE) mission. It retrieves aerosol and
surface properties from multi-angle polarimetric observations acquired
by the PACE polarimeters, HARP2 and SPEXone.

FastMAPOL combines physics-based vector radiative transfer with neural
network (NN) emulators to achieve both retrieval accuracy and
computational efficiency, enabling operational-scale processing of PACE
polarimetric observations. The algorithm jointly retrieves aerosol and
surface properties and provides associated pixel-level uncertainty
estimates.

Major features of FastMAPOL include:

\begin{itemize}
\tightlist
\item
  Coupled atmosphere-surface vector radiative transfer using the PACE
  Simulator
\item
  Cascaded neural network forward-model emulators with analytical
  Jacobians
\item
  Nonlinear least-squares optimization for joint aerosol and surface
  retrieval
\item
  Adaptive multi-angle data screening
\item
  Pixel-level retrieval uncertainty estimation
\item
  Treatment of spectral and angular uncertainty correlations
\item
  Atmospheric and bidirectional reflectance correction for angular
  surface reflectance
\end{itemize}

FastMAPOL retrieval products include aerosol optical and microphysical
properties, aerosol layer height, ocean bio-optical properties,
multi-angle remote sensing reflectance over ocean, and surface
reflectance properties over land.

This Algorithm Theoretical Basis Document (ATBD) describes the physical
basis, radiative transfer models, neural network forward models,
inversion methodology, uncertainty characterization, data products, and
validation of FastMAPOL for PACE HARP2 and SPEXone observations.

The remainder of this document is organized into chapters describing the
FastMAPOL algorithm, forward models, retrieval methodology, data
products, and validation results. A high-level overview is first
provided in algorithm summary for readers who want a concise description
of the algorithm and products before proceeding to the detailed
technical sections.

\bookmarksetup{startatroot}

\chapter{Document Information}\label{document-information}

\begin{longtable}[]{@{}
  >{\raggedright\arraybackslash}p{(\linewidth - 2\tabcolsep) * \real{0.5000}}
  >{\raggedright\arraybackslash}p{(\linewidth - 2\tabcolsep) * \real{0.5000}}@{}}
\toprule\noalign{}
\begin{minipage}[b]{\linewidth}\raggedright
Item
\end{minipage} & \begin{minipage}[b]{\linewidth}\raggedright
Description
\end{minipage} \\
\midrule\noalign{}
\endhead
\bottomrule\noalign{}
\endlastfoot
Document Title & FastMAPOL Algorithm Theoretical Basis Document \\
Version & v2026.8.26 \\
Date & 2026-08-26 \\
Authors & Meng Gao, Pengwang Zhai, Kamal Aryal, Kirk Knobelspiesse,
Bryan Franz \\
Institution & NASA Goddard Space Flight Center (GSFC) \\
Project & PACE Mission \\
\end{longtable}

\bookmarksetup{startatroot}

\chapter{Algorithm Overview}\label{algorithm-overview}

Before diving into the full document, this chapter provides a high-level
overview of the FastMAPOL algorithm, data products, product format, and
validation results.

FastMAPOL is a physics-based retrieval algorithm designed to retrieve
aerosol, ocean color, and land surface properties from multi-angle
polarimetric observations. It combines a coupled vector radiative
transfer framework with neural network (NN) forward models and numerical
optimization to achieve both physical consistency and computational
efficiency.

At a high level, FastMAPOL works through the following steps:

\begin{enumerate}
\def\labelenumi{\arabic{enumi}.}
\item
  \textbf{Forward modeling:} A coupled atmosphere-ocean or
  atmosphere-land vector radiative transfer model describes the measured
  multi-angle radiance and polarization as functions of aerosol
  properties, surface properties, and observation geometry.
\item
  \textbf{Neural network acceleration:} NNs trained using rigorous
  radiative transfer simulations serve as fast forward-model emulators,
  allowing radiance and polarization to be calculated efficiently during
  iterative retrievals.
\item
  \textbf{Joint optimization:} Aerosol and surface parameters are
  retrieved simultaneously by minimizing the differences between the
  measured and modeled multi-angle reflectance and degree of linear
  polarization (DoLP).
\item
  \textbf{Adaptive data screening:} Measurements that cannot be
  adequately represented by the forward model, for example because of
  cloud contamination or other subpixel heterogeneity, can be identified
  from the fitting residuals and adaptively excluded from subsequent
  retrieval iterations.
\item
  \textbf{Uncertainty estimation:} Pixel-level retrieval uncertainties
  are estimated using the measurement uncertainties, forward-model
  uncertainties, and Jacobian matrix, with uncertainty correlations
  considered where applicable.
\item
  \textbf{Surface reflectance retrieval:} Following the aerosol
  retrieval and atmospheric correction, FastMAPOL derives surface
  reflectance products, including angular remote sensing reflectance
  (\(R_{rs}\)) over ocean and angular surface reflectance over land.
  Bidirectional reflectance corrections can further be applied to
  provide standardized surface reflectance products.
\end{enumerate}

The resulting framework uses the information contained in multi-angle
measurements to jointly constrain atmospheric and surface properties
while maintaining consistency between aerosol retrieval, atmospheric
correction, and surface reflectance retrieval.

For more technical details, please refer to the subsequent sections,
which describe:

\begin{itemize}
\tightlist
\item
  Vector radiative transfer, including aerosol, ocean, and land surface
  models
\item
  Neural network (NN) forward model training
\item
  Analytical NN Jacobian calculation using automatic differentiation
\item
  Multi-angle adaptive data screening
\item
  Pixel-level uncertainty estimation
\item
  Treatment of uncertainty correlations
\item
  Angular surface reflectance for ocean and land
\end{itemize}

\bookmarksetup{startatroot}

\chapter{Retrieval algorithm summary}\label{retrieval-algorithm-summary}

FastMAPOL is an efficient coupled aerosol and ocean color retrieval
algorithm that combines neural network (NN) forward models with
iterative optimization (M. Gao et al. 2021; M. Gao, Franz, et al. 2023).
The NN forward models emulate vector radiative transfer calculations for
a coupled atmosphere-ocean system based on the PACE simulator (Zhai et
al. 2009, 2010; Zhai and Hu 2022). Please find more details in the NN,
automatic differentiaton and aerosol and surface model details in the
technical section.

\section{Radiative transfer
parameterization}\label{radiative-transfer-parameterization}

The radiative transfer (RT) model used in this work is the PACE
simulator (Zhai and Hu 2022), which solves the polarized radiance field
in a coupled atmosphere-ocean system with flexible solar and viewing
geometries, aerosol microphysical properties, and ocean bio-optical
models (Meng Gao et al. 2018). The RT model is used to generate the
training datasets for the neural network forward models (M. Gao et al.
2021; M. Gao, Franz, et al. 2023).

The forward radiative transfer simulations use atmospheric molecular
vertical distributions from the U.S. Standard Atmosphere profile
(Anderson et al. 1986), scaled by surface pressure. Absorption by
oxygen, water vapor, methane, carbon dioxide, ozone, and nitrogen
dioxide is accounted for using line-by-line calculations and the
double-k distribution method (Duan, Min, and Li 2005; Zhai and Hu 2022).
Surface pressure \(P_s\) and ozone column density \(n_{\mathrm{O}_3}\)
are considered as input parameters, with values obtained from the
MERRA-2 reanalysis (Gelaro et al. 2017).

Aerosols are represented by a multi-modal framework that allows
independent characterization of fine- and coarse-mode size
distributions, each with its own complex refractive index defined by
real (\(m_r\)) and imaginary (\(m_i\)) components, assuming flat
spectral dependence. The fine mode is represented by three submodes,
while the coarse mode is represented by two submodes. Each submode is
described by a log-normal size distribution with mean radii of 0.1,
0.1732, 0.3, 1.0, and 2.9 \(\mu\mathrm{m}\) and geometric standard
deviations of 0.35, 0.35, 0.35, 0.5, and 0.5, respectively, consistent
with Oleg Dubovik et al. (2006) and Xu et al. (2016). The volume density
of each submode is denoted by \(V_i\).

Both fine and coarse modes allow mixtures of spherical and spheroidal
particles to account for particle non-sphericity (Oleg Dubovik et al.
2006). The corresponding spherical fractions are represented by
\(\mathrm{sph}_f\) and \(\mathrm{sph}_c\). The aerosol vertical
distribution is parameterized by the aerosol layer height (\(z_c\)),
defined as the center of a Gaussian vertical profile describing aerosol
number density.

Ocean bio-optical properties, including contributions from pure
seawater, phytoplankton, and colored dissolved organic matter (CDOM),
are parameterized using chlorophyll-a concentration (Chl-a). The
scattering phase functions of phytoplankton are represented by the
Fournier-Forand (FF) phase function (Fournier and Forand 1994), which is
mixed with the pure seawater phase function to describe total scattering
in seawater. The Mueller matrix is represented by the product of the
mixed phase function and the average normalized Mueller matrix measured
by Voss and Fry (1984) and parameterized by Kokhanovsky (2003).

Ocean surface roughness is described using the Cox-Munk model,
parameterized by surface wind speed (\(w_s\)). The whitecap fraction is
parameterized in terms of wind speed, with whitecap reflectance based on
Koepke (1984).

The current bio-optical model is best suited for waters optically
dominated by phytoplankton and performs well over most global ocean
waters (Meng Gao et al. 2026), but it faces challenges in accurately
representing optically complex coastal waters, particularly those
influenced by non-algal particles and strong CDOM absorption. Ongoing
development of more advanced bio-optical models capable of representing
complex water types will enable more realistic retrievals in coastal
waters, particularly in regions influenced by both complex aerosol
properties and bright ocean constituents Aryal et al. (2026).

All forward-model parameters and their ranges are summarized in
Table~\ref{tbl-param-ranges}. Solar and viewing geometries are specified
by the solar zenith angle (\(\theta_0\)), viewing zenith angle
(\(\theta_v\)), and relative azimuth angle (\(\phi_v\)), which are
obtained from PACE Level-1C data. Surface pressure and ozone are
obtained from ancillary data, while the aerosol, ocean, and surface
parameters are retrieved by FastMAPOL.

\begin{longtable}[]{@{}lrrr@{}}
\caption{Parameters and ranges used in the neural network forward model.
The minimum and maximum values define the bounds used to sample the NN
training dataset and serve as constraints in the retrieval
algorithm.}\label{tbl-param-ranges}\tabularnewline
\toprule\noalign{}
\textbf{Parameter} & \textbf{Unit} & \textbf{Min} & \textbf{Max} \\
\midrule\noalign{}
\endfirsthead
\toprule\noalign{}
\textbf{Parameter} & \textbf{Unit} & \textbf{Min} & \textbf{Max} \\
\midrule\noalign{}
\endhead
\bottomrule\noalign{}
\endlastfoot
\(\theta_0\) & \(^\circ\) & 0 & 80 \\
\(\theta_v\) & \(^\circ\) & 0 & 80 \\
\(\phi_v\) & \(^\circ\) & 0 & 180 \\
\(n_{\mathrm{O}_3}\) & Dobson & 150 & 450 \\
\(P_s\) & mb & 950 & 1050 \\
\(m_{r,f}\) & -- & 1.33 & 1.55 \\
\(m_{r,c}\) & -- & 1.33 & 1.55 \\
\(m_{i,f}\) & -- & 0 & 0.03 \\
\(m_{i,c}\) & -- & 0 & 0.001 \\
\(V_1\) & \(\mu\mathrm{m}^3\,\mu\mathrm{m}^{-2}\) & 0 & 0.14 \\
\(V_2\) & \(\mu\mathrm{m}^3\,\mu\mathrm{m}^{-2}\) & 0 & 0.11 \\
\(V_3\) & \(\mu\mathrm{m}^3\,\mu\mathrm{m}^{-2}\) & 0 & 0.07 \\
\(V_4\) & \(\mu\mathrm{m}^3\,\mu\mathrm{m}^{-2}\) & 0 & 0.20 \\
\(V_5\) & \(\mu\mathrm{m}^3\,\mu\mathrm{m}^{-2}\) & 0 & 0.62 \\
\(\mathrm{sph}_f\) & -- & 0 & 1 \\
\(\mathrm{sph}_c\) & -- & 0 & 1 \\
\(z_c\) & km & 0.1 & 6.0 \\
\(w_s\) & \(\mathrm{m}\,\mathrm{s}^{-1}\) & 0.5 & 15 \\
Chl-a & \(\mathrm{mg}\,\mathrm{m}^{-3}\) & 0.01 & 10 \\
\end{longtable}

Note that aerosol retrieval over land has also been implemented for
HARP2, with the corresponding product available as MAPOL\_LAND. The
retrieval follows the same coupled atmosphere-surface radiative transfer
framework. Additional technical details are provided in later sections.

\section{Neural network forward
models}\label{neural-network-forward-models}

Neural network models are employed to represent TOA reflectance
(\(\rho_t\)) and degree of linear polarization (DoLP) with two levels of
accuracy. A smaller, computationally efficient neural network is used to
generate an initial first guess, while a larger and more accurate neural
network is applied in subsequent refinement steps of the retrieval. This
two-stage strategy balances computational efficiency with retrieval
accuracy (M. Gao, Franz, et al. 2023).

In addition to the TOA forward models, separate neural network models
are developed to support atmospheric correction. These include NNs
trained to represent atmospheric and surface path reflectance
(\(\rho_{\mathrm{atm+surface}}\)), total transmittance
(\(1/(T_d t_u)\)), and a combined transmittance and BRDF correction term
(\(C_{\mathrm{BRDF}}/(T_d t_u)\)). These components enable efficient
decomposition of the TOA signal and facilitate the retrieval of
water-leaving reflectance.

Additional neural network models are trained to compute AOD and SSA for
both fine and coarse modes, allowing consistent treatment of aerosol
optical and microphysical properties within the coupled retrieval
framework. All NN models, architectures, and associated uncertainties
are summarized in Table~\ref{tbl-nn-arch}.

\begin{longtable}[]{@{}
  >{\raggedright\arraybackslash}p{(\linewidth - 4\tabcolsep) * \real{0.3000}}
  >{\raggedright\arraybackslash}p{(\linewidth - 4\tabcolsep) * \real{0.3000}}
  >{\raggedleft\arraybackslash}p{(\linewidth - 4\tabcolsep) * \real{0.4000}}@{}}
\caption{Neural network architectures and associated uncertainties used
in FastMAPOL for SPEXone (34 spectral bands). For HARP2 the output bands
are 4 with similar structures (M. Gao, Franz, et al.
2023).}\label{tbl-nn-arch}\tabularnewline
\toprule\noalign{}
\begin{minipage}[b]{\linewidth}\raggedright
\textbf{Quantity}
\end{minipage} & \begin{minipage}[b]{\linewidth}\raggedright
\textbf{Neural network architecture}
\end{minipage} & \begin{minipage}[b]{\linewidth}\raggedleft
\textbf{Uncertainty}
\end{minipage} \\
\midrule\noalign{}
\endfirsthead
\toprule\noalign{}
\begin{minipage}[b]{\linewidth}\raggedright
\textbf{Quantity}
\end{minipage} & \begin{minipage}[b]{\linewidth}\raggedright
\textbf{Neural network architecture}
\end{minipage} & \begin{minipage}[b]{\linewidth}\raggedleft
\textbf{Uncertainty}
\end{minipage} \\
\midrule\noalign{}
\endhead
\bottomrule\noalign{}
\endlastfoot
\(\rho_t\) & \(19\times256\times256\times256\times34\) & 0.5\% \\
\(\rho_t\) (first guess) & \(19\times128\times128\times128\times34\) &
0.6\% \\
\(\rho_{\mathrm{atm+surface}}\) &
\(18\times256\times256\times256\times34\) & 0.5\% \\
DoLP & \(19\times512\times512\times512\times34\) & 0.0007 \\
DoLP (first guess) & \(19\times256\times256\times256\times34\) &
0.0012 \\
\(C_{\mathrm{BRDF}}/(T_d t_u)\) & \(19\times256\times256\times34\) &
0.006 \\
\(1/(T_d t_u)\) & \(19\times256\times256\times34\) & 0.005 \\
AOD (fine) & \(6\times64\times64\times34\) & 0.003 \\
AOD (coarse) & \(5\times64\times64\times34\) & 0.001 \\
SSA (fine) & \(6\times64\times64\times34\) & 0.003 \\
SSA (coarse) & \(5\times64\times64\times34\) & 0.01 \\
\end{longtable}

\section{Retrieval and optimization}\label{retrieval-and-optimization}

FastMAPOL determines the optimal values of the state parameters by
minimizing the difference between the measurements and NN forward model
predictions. An iterative optimization approach is used. The NN forward
model computes the reflectance and degree of linear polarization (DoLP),
defined as

\begin{equation}\phantomsection\label{eq-rho}{
\rho_t
=
\frac{\pi L_t}{\mu_0 F_0}
}\end{equation}

and

\begin{equation}\phantomsection\label{eq-DoLP}{
P_t
=
\frac{\sqrt{Q_t^2+U_t^2}}{L_t},
}\end{equation}

where \(L_t\), \(Q_t\), and \(U_t\) are the Stokes parameters at the
sensor altitude, \(F_0\) is the extraterrestrial solar irradiance, and
\(\mu_0\) is the cosine of the solar zenith angle.

The differences between the measurements and model predictions are
represented by a cost function, \(\chi^2\), based on Bayesian theory
(Rodgers 2000):

\begin{equation}\phantomsection\label{eq-cost}{
\chi^2({\mathbf x})
=
\frac{1}{N}
\sum_i
\left(
\frac{[\rho_t(i)-\rho_t^f({\mathbf x};i)]^2}{\sigma_\rho^2(i)}
+
\frac{[P_t(i)-P_t^f({\mathbf x};i)]^2}{\sigma_P^2(i)}
\right).
}\end{equation}

Here, \(\rho_t^f\) and \(P_t^f\) are the reflectance and DoLP computed
by the NN forward model. The state vector \({\mathbf x}\) contains the
retrieved aerosol, ocean, and surface parameters described in
Table~\ref{tbl-param-ranges}. The subscript \(i\) denotes the
measurement index, where one measurement is defined as a pair of
reflectance and DoLP values at a given viewing angle and wavelength, and
\(N\) is the total number of measurements used in the retrieval.

The total uncertainties of reflectance and DoLP used in the cost
function are denoted by \(\sigma_\rho\) and \(\sigma_P\), respectively.
These include contributions from instrument uncertainties
\(\sigma_{\mathrm{ins}}\), NN forward model uncertainties
\(\sigma_{\mathrm{NN}}\), and radiative transfer simulation
uncertainties \(\sigma_{\mathrm{RT}}\) used to train the NN:

\begin{equation}\phantomsection\label{eq-sigma-rho}{
\sigma_\rho^2
=
\sigma_{\rho,\mathrm{ins}}^2
+
\sigma_{\rho,\mathrm{NN}}^2
+
\sigma_{\rho,\mathrm{RT}}^2,
}\end{equation}

\begin{equation}\phantomsection\label{eq-sigma-P}{
\sigma_P^2
=
\sigma_{P,\mathrm{ins}}^2
+
\sigma_{P,\mathrm{NN}}^2
+
\sigma_{P,\mathrm{RT}}^2.
}\end{equation}

The NN forward-model uncertainties are determined during NN training and
evaluation (M. Gao et al. 2021; M. Gao, Franz, et al. 2023), with the
current NN architectures and associated uncertainties summarized in
Table~\ref{tbl-nn-arch}. At present, uncertainties are assumed to be
spectrally and angularly uncorrelated.

During FastMAPOL retrievals, with each iteration, the state vector
\({\mathbf x}\) is updated, and the NN forward model is used to compute
\(\rho_t^f\) and \(P_t^f\), from which a new value of the cost function
is calculated. The Jacobian matrix is used to determine the magnitude
and direction of the update to \({\mathbf x}\). The NN forward model and
analytical calculation of the Jacobian matrix through automatic
differentiation are described in the following section.

This iterative process continues until the convergence criterion is met:

\begin{equation}\phantomsection\label{eq-delta-chi}{
\frac{|\chi_i^2-\chi_{i-1}^2|}
{\chi_i^2}
<
\eta,
}\end{equation}

where \(i\) is the iteration index and \(\eta\) is the convergence
tolerance, taken as 0.01.

The convergence and overall quality of the retrieval can be evaluated
using the cost function. For an ensemble of retrievals, if the residuals
are consistent with the assumed uncertainties, the normalized
cost-function distribution can be compared with the theoretical
probability density function (PDF) of the \(\chi^2\) distribution:

\begin{equation}\phantomsection\label{eq-chi2-pdf}{
f(\chi^2,k)
=
\frac{(\chi^2)^{k/2-1} k^{k/2} e^{-\chi^2 k/2}}
{2^{k/2}\Gamma(k/2)}.
}\end{equation}

Here, \(\chi^2\) is the cost function defined in Equation~\ref{eq-cost},
\(k\) is the number of degrees of freedom (DOF), and \(\Gamma(k/2)\)
denotes the gamma function (James 2006). Neglecting potential
correlations among measurement uncertainties, the DOF can be
approximated by the total number of measurements, \(N\). More
rigorously, the effective DOF depends on the number and information
content of the retrieved parameters and can be evaluated using the
Jacobian and error covariance matrices (Rodgers 2000).

\bookmarksetup{startatroot}

\chapter{Data Products and Format}\label{data-products-and-format}

\section{Overview}\label{overview}

FastMAPOL (Fast Multi-Angle Polarimetric Ocean and Land algorithm)
produces Level-2 (L2) geophysical retrievals of aerosol and surface
properties from the multi-angle polarimetric measurements of the
\textbf{PACE polarimeters}, including \textbf{HARP2} and
\textbf{SPEXone}.

The Level-2 product provides pixel-level estimates of aerosol optical
properties, aerosol microphysical parameters, and ocean color quantities
derived through atmospheric correction. In addition to the primary
geophysical retrievals, the product includes diagnostic quantities
describing retrieval performance and quality indicators for downstream
processing.

The Level-2 product is distributed in \textbf{NetCDF format} and
organized into hierarchical data groups to separate geolocation,
geophysical variables, diagnostics, and sensor parameters.

The FastMAPOL Level-2 product is designed to support:

\begin{itemize}
\tightlist
\item
  aerosol science applications
\item
  ocean color retrievals
\item
  atmospheric correction
\item
  algorithm validation
\item
  Level-3 gridded product generation
\end{itemize}

\begin{center}\rule{0.5\linewidth}{0.5pt}\end{center}

\section{File Organization}\label{file-organization}

The FastMAPOL Level-2 file is organized into several top-level groups.
Each group contains variables associated with a specific component of
the retrieval system.

\begin{longtable}[]{@{}ll@{}}
\toprule\noalign{}
Group & Description \\
\midrule\noalign{}
\endhead
\bottomrule\noalign{}
\endlastfoot
geolocation\_data & Pixel geolocation information \\
geophysical\_data & Primary aerosol and ocean retrieval products \\
diagnostic\_data & Retrieval diagnostics and performance metrics \\
sensor\_band\_parameters & Instrument spectral and angular parameters \\
processing\_control & Metadata and algorithm configuration \\
\end{longtable}

\subsection{Dimensions}\label{dimensions}

The FastMAPOL Level-2 product uses a set of core dimensions describing
spatial, spectral, and angular sampling.

\begin{longtable}[]{@{}ll@{}}
\toprule\noalign{}
Dimension & Description \\
\midrule\noalign{}
\endhead
\bottomrule\noalign{}
\endlastfoot
number\_of\_lines & Number of along-track scan lines \\
pixels\_per\_line & Number of cross-track pixels \\
wavelength & Spectral bands used in retrieval \\
number\_of\_views & Viewing angles used in the retrieval \\
intensity\_bands\_per\_view & Intensity channels per view \\
polarization\_bands\_per\_view & Polarization channels per view \\
\end{longtable}

\subsection{Instrument Differences}\label{instrument-differences}

\textbf{HARP2}

\begin{itemize}
\tightlist
\item
  \textasciitilde90 viewing angles\\
\item
  4 spectral bands : 440, 550, 670, 870nm
\item
  wide swath (\textasciitilde2000km)
\end{itemize}

\textbf{SPEXone}

\begin{itemize}
\tightlist
\item
  5 viewing angles\\
\item
  hyperspectral polarimetric measurements 50
\item
  narrow swath (\textasciitilde100km)
\end{itemize}

A total of 34 specific wavelengths from the hyperspectral measurements
for aerosol and surface retrieval.

\section{Geolocation Data}\label{geolocation-data}

Group: \texttt{/geolocation\_data}

\begin{longtable}[]{@{}lll@{}}
\toprule\noalign{}
Variable & Dimensions & Description \\
\midrule\noalign{}
\endhead
\bottomrule\noalign{}
\endlastfoot
latitude & (number\_of\_lines, pixels\_per\_line) & Pixel latitude \\
longitude & (number\_of\_lines, pixels\_per\_line) & Pixel longitude \\
\end{longtable}

These variables define the geographic location of each retrieval pixel.

\section{Geophysical Retrieval
Products}\label{geophysical-retrieval-products}

Group: \texttt{/geophysical\_data}

The geophysical group contains aerosol optical properties, aerosol
microphysical parameters, and ocean color products retrieved by
FastMAPOL. Two data suites, MAPOL\_OCEAN and MAPOL\_LAND, are available
for aerosol over ocean and land, respectively.

\subsection{Aerosol Optical Properties (MAPOL\_OCEAN and
MAPOL\_LAND)}\label{aerosol-optical-properties-mapol_ocean-and-mapol_land}

\begin{longtable}[]{@{}
  >{\raggedright\arraybackslash}p{(\linewidth - 4\tabcolsep) * \real{0.2778}}
  >{\raggedright\arraybackslash}p{(\linewidth - 4\tabcolsep) * \real{0.3611}}
  >{\raggedright\arraybackslash}p{(\linewidth - 4\tabcolsep) * \real{0.3611}}@{}}
\toprule\noalign{}
\begin{minipage}[b]{\linewidth}\raggedright
Variable
\end{minipage} & \begin{minipage}[b]{\linewidth}\raggedright
Dimensions
\end{minipage} & \begin{minipage}[b]{\linewidth}\raggedright
Description
\end{minipage} \\
\midrule\noalign{}
\endhead
\bottomrule\noalign{}
\endlastfoot
aot & (number\_of\_lines, pixels\_per\_line, wavelength) & Aerosol
optical depth \\
aot\_fine & (number\_of\_lines, pixels\_per\_line, wavelength) &
Fine-mode optical depth \\
aot\_coarse & (number\_of\_lines, pixels\_per\_line, wavelength) &
Coarse-mode optical depth \\
ssa & (number\_of\_lines, pixels\_per\_line, wavelength) & Aerosol
single scattering albedo \\
ssa\_fine & (number\_of\_lines, pixels\_per\_line, wavelength) &
Fine-mode SSA \\
ssa\_coarse & (number\_of\_lines, pixels\_per\_line, wavelength) &
Coarse-mode SSA \\
angstrom\_440\_870 & (number\_of\_lines, pixels\_per\_line) & Ångström
exponent (440-870 nm) \\
angstrom\_440\_670 & (number\_of\_lines, pixels\_per\_line) & Ångström
exponent (440-670 nm) \\
fmf & (number\_of\_lines, pixels\_per\_line) & Fine-mode optical depth
fraction \\
\end{longtable}

These parameters describe aerosol optical loading and spectral behavior.
Note that Ångström exponent (440-870 nm) is only available for HARP2
data since SPEXone do not have 870nm band.

\subsection{Aerosol Microphysical Properties (MAPOL\_OCEAN and
MAPOL\_LAND)}\label{aerosol-microphysical-properties-mapol_ocean-and-mapol_land}

\begin{longtable}[]{@{}
  >{\raggedright\arraybackslash}p{(\linewidth - 4\tabcolsep) * \real{0.2778}}
  >{\raggedright\arraybackslash}p{(\linewidth - 4\tabcolsep) * \real{0.3611}}
  >{\raggedright\arraybackslash}p{(\linewidth - 4\tabcolsep) * \real{0.3611}}@{}}
\toprule\noalign{}
\begin{minipage}[b]{\linewidth}\raggedright
Variable
\end{minipage} & \begin{minipage}[b]{\linewidth}\raggedright
Dimensions
\end{minipage} & \begin{minipage}[b]{\linewidth}\raggedright
Description
\end{minipage} \\
\midrule\noalign{}
\endhead
\bottomrule\noalign{}
\endlastfoot
reff\_fine & (number\_of\_lines, pixels\_per\_line) & Fine-mode
effective radius \\
reff\_coarse & (number\_of\_lines, pixels\_per\_line) & Coarse-mode
effective radius \\
veff\_fine & (number\_of\_lines, pixels\_per\_line) & Fine-mode
effective variance \\
veff\_coarse & (number\_of\_lines, pixels\_per\_line) & Coarse-mode
effective variance \\
mr\_fine & (number\_of\_lines, pixels\_per\_line, wavelength) &
Fine-mode Real refractive index \\
mr\_coarse & (number\_of\_lines, pixels\_per\_line, wavelength) &
Coarse-mode Real refractive index \\
mr & (number\_of\_lines, pixels\_per\_line, wavelength) & Mode-average
Real refractive index \\
mi\_fine & (number\_of\_lines, pixels\_per\_line, wavelength) &
Fine-mode Imaginary refractive index \\
mi\_coarse & (number\_of\_lines, pixels\_per\_line, wavelength) &
Coarse-mode Imaginary refractive index \\
mi & (number\_of\_lines, pixels\_per\_line, wavelength) & Mode-average
Imaginary refractive index \\
sph\_fine & (number\_of\_lines, pixels\_per\_line) & Fine-mode Spherical
particle fraction \\
sph\_coarse & (number\_of\_lines, pixels\_per\_line) & Coarse-mode
Spherical particle fraction \\
sph & (number\_of\_lines, pixels\_per\_line) & Volum-average Spherical
particle fraction \\
vd\_fine & (number\_of\_lines, pixels\_per\_line) & Fine-mode volume
density \\
vd\_coarse & (number\_of\_lines, pixels\_per\_line) & Coarse-mode volume
density \\
fvf & (number\_of\_lines, pixels\_per\_line) & Fine-mode volume
fraction \\
alh & (number\_of\_lines, pixels\_per\_line) & Aerosol layer height \\
\end{longtable}

These variables constrain aerosol particle size distribution,
composition, and shape.

\subsection{Ocean Products
(MAPOL\_OCEAN)}\label{ocean-products-mapol_ocean}

\begin{longtable}[]{@{}
  >{\raggedright\arraybackslash}p{(\linewidth - 4\tabcolsep) * \real{0.2778}}
  >{\raggedright\arraybackslash}p{(\linewidth - 4\tabcolsep) * \real{0.3611}}
  >{\raggedright\arraybackslash}p{(\linewidth - 4\tabcolsep) * \real{0.3611}}@{}}
\toprule\noalign{}
\begin{minipage}[b]{\linewidth}\raggedright
Variable
\end{minipage} & \begin{minipage}[b]{\linewidth}\raggedright
Dimensions
\end{minipage} & \begin{minipage}[b]{\linewidth}\raggedright
Description
\end{minipage} \\
\midrule\noalign{}
\endhead
\bottomrule\noalign{}
\endlastfoot
Rrs\_angular & (number\_of\_lines, pixels\_per\_line, wavelength,
number\_of\_views) & Angular remote sensing reflectance before BRDF
correction \\
Rrs\_nadir & (number\_of\_lines, pixels\_per\_line, wavelength,
number\_of\_views) & Nadir adjusted remote sensing reflectance after
BRDF correction \\
Rrs\_angular\_mean & (number\_of\_lines, pixels\_per\_line, wavelength)
& Angular mean of Rrs\_angular \\
Rrs\_angular\_std & (number\_of\_lines, pixels\_per\_line, wavelength) &
Angular standard deviation of Rrs\_angular \\
Rrs\_nadir\_mean & (number\_of\_lines, pixels\_per\_line, wavelength) &
Angular mean of Rrs\_nadir \\
Rrs\_nadir\_std & (number\_of\_lines, pixels\_per\_line, wavelength) &
Angular standard deviation of Rrs\_nadir \\
wind\_speed & (number\_of\_lines, pixels\_per\_line) & Surface wind
speed \\
chla & (number\_of\_lines, pixels\_per\_line) & Chlorophyll-a
concentration \\
\end{longtable}

The remote sensing reflectance represents ocean-leaving reflectance
after atmospheric correction. To compare with AERONET OC or other
dataset for validation purpose after BRDF correction, variable
Rrs\_nadir\_mean is suggested to use.

Note that V3 data contains \textbf{Rrs1} and \textbf{Rrs2}, which are
renmaed as \textbf{Rrs\_angular} and \textbf{Rrs\_nadir} in current V4
data.

\subsection{Land surface product
(MAPOL\_LAND)}\label{land-surface-product-mapol_land}

\begin{longtable}[]{@{}
  >{\raggedright\arraybackslash}p{(\linewidth - 4\tabcolsep) * \real{0.2778}}
  >{\raggedright\arraybackslash}p{(\linewidth - 4\tabcolsep) * \real{0.3611}}
  >{\raggedright\arraybackslash}p{(\linewidth - 4\tabcolsep) * \real{0.3611}}@{}}
\toprule\noalign{}
\begin{minipage}[b]{\linewidth}\raggedright
Variable
\end{minipage} & \begin{minipage}[b]{\linewidth}\raggedright
Dimensions
\end{minipage} & \begin{minipage}[b]{\linewidth}\raggedright
Description
\end{minipage} \\
\midrule\noalign{}
\endhead
\bottomrule\noalign{}
\endlastfoot
rhos\_angular & (number\_of\_lines, pixels\_per\_line, wavelength,
number\_of\_views) & Angular land surface reflectance before BRDF
correction \\
rhos\_nadir & (number\_of\_lines, pixels\_per\_line, wavelength,
number\_of\_views) & Nadir adjusted land surface reflectance after BRDF
correction \\
rhos\_angular\_mean & (number\_of\_lines, pixels\_per\_line, wavelength)
& Angular mean of rhos\_angular \\
rhos\_angular\_std & (number\_of\_lines, pixels\_per\_line, wavelength)
& Angular standard deviation of rhos\_angular \\
rhos\_nadir\_mean & (number\_of\_lines, pixels\_per\_line, wavelength) &
Angular mean of rhos\_nadir \\
rhos\_nadir\_std & (number\_of\_lines, pixels\_per\_line, wavelength) &
Angular standard deviation of rhos\_nadir \\
\end{longtable}

The land surface reflectance represents reflectance after atmospheric
correction. To compare with validation dataset after BRDF correction,
variable rhos\_nadir\_mean is suggested to use.

\section{Diagnostic Data}\label{diagnostic-data}

Group: \texttt{/diagnostic\_data}

These variables describe retrieval convergence and performance.

\begin{longtable}[]{@{}
  >{\raggedright\arraybackslash}p{(\linewidth - 4\tabcolsep) * \real{0.2778}}
  >{\raggedright\arraybackslash}p{(\linewidth - 4\tabcolsep) * \real{0.3611}}
  >{\raggedright\arraybackslash}p{(\linewidth - 4\tabcolsep) * \real{0.3611}}@{}}
\toprule\noalign{}
\begin{minipage}[b]{\linewidth}\raggedright
Variable
\end{minipage} & \begin{minipage}[b]{\linewidth}\raggedright
Dimensions
\end{minipage} & \begin{minipage}[b]{\linewidth}\raggedright
Description
\end{minipage} \\
\midrule\noalign{}
\endhead
\bottomrule\noalign{}
\endlastfoot
chi2 & (number\_of\_lines, pixels\_per\_line) & Retrieval cost
function \\
chi2\_first\_guess & (number\_of\_lines, pixels\_per\_line) & Cost
function at first guess \\
nv\_ref & (number\_of\_lines, pixels\_per\_line) & Number of reflectance
measurements used \\
nv\_dolp & (number\_of\_lines, pixels\_per\_line) & Number of
polarization measurements used \\
nfev & (number\_of\_lines, pixels\_per\_line) & Number of forward model
evaluations \\
njev & (number\_of\_lines, pixels\_per\_line) & Number of Jacobian
evaluations \\
timing & (number\_of\_lines, pixels\_per\_line) & Retrieval runtime \\
quality\_flag & (number\_of\_lines, pixels\_per\_line) & Retrieval
quality indicator \\
ozone & (number\_of\_lines, pixels\_per\_line) & Total column ozone \\
surface\_pressure & (number\_of\_lines, pixels\_per\_line) & Surface
pressure \\
\end{longtable}

\subsection{Retrieval Quality Metrics}\label{retrieval-quality-metrics}

FastMAPOL retrievals are performed using an optimal estimation framework
that minimizes the cost function

\[
\chi^2 = \frac{1}{N}\sum \frac{(f-m)^2}{\sigma^2}
\]

where

\begin{itemize}
\tightlist
\item
  \(m\) is the measurement\\
\item
  \(f\) is the forward model simulation
\item
  \(\sigma\) is measurement uncertainty\\
\item
  \(N\) is the number of measurements used in the retrieval.
\end{itemize}

Lower values of \(\chi^2\) indicate better agreement between
measurements and forward model simulations.

\subsection{Retrieval Quality Flag}\label{retrieval-quality-flag}

A quality flag is provided to facilitate filtering and Level-3
processing.

\begin{longtable}[]{@{}ll@{}}
\toprule\noalign{}
quality\_flag & Description \\
\midrule\noalign{}
\endhead
\bottomrule\noalign{}
\endlastfoot
0 & highest retrieval quality \\
1 & good retrieval quality \\
\textgreater1 & reduced quality \\
\end{longtable}

The quality flag is determined from combinations of:

\begin{itemize}
\tightlist
\item
  cost function value
\item
  number of reflectance measurements
\item
  number of polarization measurements.
\end{itemize}

\subsection{Measurement Screening
Masks}\label{measurement-screening-masks}

FastMAPOL performs adaptive measurement screening during the inversion.

\begin{longtable}[]{@{}
  >{\raggedright\arraybackslash}p{(\linewidth - 4\tabcolsep) * \real{0.2778}}
  >{\raggedright\arraybackslash}p{(\linewidth - 4\tabcolsep) * \real{0.3611}}
  >{\raggedright\arraybackslash}p{(\linewidth - 4\tabcolsep) * \real{0.3611}}@{}}
\toprule\noalign{}
\begin{minipage}[b]{\linewidth}\raggedright
Variable
\end{minipage} & \begin{minipage}[b]{\linewidth}\raggedright
Dimensions
\end{minipage} & \begin{minipage}[b]{\linewidth}\raggedright
Description
\end{minipage} \\
\midrule\noalign{}
\endhead
\bottomrule\noalign{}
\endlastfoot
mask\_ref & (number\_of\_lines, pixels\_per\_line, number\_of\_views,
intensity\_bands\_per\_view) & Reflectance measurement mask \\
mask\_dolp & (number\_of\_lines, pixels\_per\_line, number\_of\_views,
polarization\_bands\_per\_view) & Polarization measurement mask \\
\end{longtable}

Mask interpretation:

\begin{longtable}[]{@{}ll@{}}
\toprule\noalign{}
Value & Meaning \\
\midrule\noalign{}
\endhead
\bottomrule\noalign{}
\endlastfoot
0 & measurement used in retrieval \\
1 or NaN & measurement excluded from retrieval \\
\end{longtable}

These masks allow users to evaluate measurement quality at each viewing
angle.

\section{Sensor Band Parameters}\label{sensor-band-parameters}

Group: \texttt{/sensor\_band\_parameters}

\begin{longtable}[]{@{}lll@{}}
\toprule\noalign{}
Variable & Dimensions & Description \\
\midrule\noalign{}
\endhead
\bottomrule\noalign{}
\endlastfoot
wavelength & (wavelength) & Retrieval wavelengths \\
sensor\_view\_angle & (number\_of\_views) & Viewing angle geometry \\
\end{longtable}

\bookmarksetup{startatroot}

\chapter{Jupyter notebook tutorials}\label{jupyter-notebook-tutorials}

Notebook tutorials are available through the
\href{https://nasa.github.io/oceandata-notebooks/sections/cloud-atmosphere.html}{PACE
Help Hub}.

\begin{itemize}
\tightlist
\item
  \href{https://nasa.github.io/oceandata-notebooks/notebooks/harp2/harp2_l2_fastmapol_aerosol_ocean_product.html}{HARP2
  Ocean Aerosols (FastMAPOL)} - Working with aerosol data over ocean
  from the HARP2 polarimeter with the FastMAPOL algorithm
\item
  \href{https://nasa.github.io/oceandata-notebooks/notebooks/harp2/harp2_l2_fastmapol_aerosol_land_product.html}{HARP2
  Land Aerosols (FastMAPOL)} - Working with aerosol data over land from
  the HARP2 polarimeter with the FastMAPOL algorithm
\item
  \href{https://nasa.github.io/oceandata-notebooks/notebooks/spexone/spexone_l2_fastmapol_aerosol_product.html}{SPEXone
  Aerosols (FastMAPOL)} - Analyzing aerosol retrievals over ocean from
  SPEXone measurements with the FastMAPOL algorithm
\end{itemize}

Please contact Meng Gao (meng.gao@nasa.gov) for additional information
or tutorial content on FastMAPOL products.

\bookmarksetup{startatroot}

\chapter{SPEXone aerosol and ocean product
validation}\label{spexone-aerosol-and-ocean-product-validation}

This page summarizes validation results for aerosol and ocean products
retrieved from the PACE multi-angle polarimeters using the
\textbf{FastMAPOL} algorithm, with emphasis on \textbf{SPEXone} aerosol
optical properties and ocean remote-sensing reflectance (\(R_{rs}\)).

FastMAPOL performs coupled aerosol--surface retrievals using an
optimal-estimation framework accelerated by deep neural network forward
models. This page focuses only on the validation results and their
interpretation.

\textbf{Additional validation results are available in Meng Gao et al.
(2026) and will be updated on the
\href{https://pace.oceansciences.org/pace_fastmapol.htm}{PACE FastMAPOL
validation page}.}

\section{Summary Information}\label{summary-information}

\begin{itemize}
\tightlist
\item
  \textbf{FastMAPOL version:} v3.0\\
\item
  \textbf{Data period:} 2024/03--2025/12\\
\item
  \textbf{Reference:} AERONET v3, Level 1.5
\end{itemize}

\section{Retrieval Products
Evaluated}\label{retrieval-products-evaluated}

\textbf{Aerosol} - Total, fine-mode, and coarse-mode aerosol optical
depth (AOD) - Total, fine-mode, and coarse-mode single-scattering albedo
(SSA)

\textbf{Ocean} - Ocean remote-sensing reflectance (\(R_{rs}\))

\section{Validation Strategy}\label{validation-strategy}

Validation is organized into three levels:

\begin{enumerate}
\def\labelenumi{\arabic{enumi}.}
\item
  \textbf{Global Level-3 assessment}\\
  Evaluate large-scale spatial and seasonal patterns as shown below.
\item
  \textbf{Global aerosol validation}\\
  Compare retrievals with AERONET observations
\item
  \textbf{Joint aerosol--ocean validation}\\
  Evaluate aerosol and \(R_{rs}\) simultaneously using AERONET-OC
  matchups, study the mutual dependencies on the retrieval
  uncertainties.
\end{enumerate}

The objective is to assess both the physical realism of the retrieved
fields and the consistency of aerosol and ocean products.

\section{Global Level-3 Products}\label{global-level-3-products}

\textbf{Aerosol Products (556 nm)}

Global distributions of total, fine-mode, and coarse-mode AOD and SSA at
556 nm:

\begin{figure}[H]

{\centering \pandocbounded{\includegraphics[keepaspectratio]{../chapters/../figure/validation_spexone_fig_15_0.04/fig_l3_aerosol_556.png}}

}

\caption{Global aerosol products}

\end{figure}%

\begin{center}\rule{0.5\linewidth}{0.5pt}\end{center}

\textbf{Ocean Products}

Global distributions of \(R_{rs}\) (443, 556, 667 nm) and
chlorophyll-\emph{a}:

\begin{figure}[H]

{\centering \pandocbounded{\includegraphics[keepaspectratio]{../chapters/../figure/validation_spexone_fig_15_0.04/fig_l3_rrs.png}}

}

\caption{Global ocean products}

\end{figure}%

The retrieved \(R_{rs}\) contains spectral and angular dimensions. The
results shown here represent the BRDF-corrected angular mean.

\textbf{Key result:}\\
The global distributions show realistic spatial and seasonal patterns,
indicating effective separation of atmospheric and oceanic retrieval
components.

\bookmarksetup{startatroot}

\chapter{HARP2 L2 aerosol product
validation}\label{harp2-l2-aerosol-product-validation}

\section{Overview}\label{overview-1}

Pre-launch Day-in-the-Life (DITL) simulations provide an end-to-end
benchmark of HARP2 measurement realism, L2 retrieval performance, and
uncertainty characterization. The results below highlight three key
aspects: synthetic observations, pixel-wise retrieval performance, and
uncertainty analysis. Details are provided in M. Gao, Franz, et al.
(2023).

Validation of the production data (V3 and V4) is ongoing, following the
methodologies reported in Meng Gao et al. (2026). Updated validation
results will be provided on the
\href{https://pace.oceansciences.org/pace_fastmapol.htm}{PACE FastMAPOL
validation page}.

\section{Summary Information}\label{summary-information-1}

\begin{itemize}
\tightlist
\item
  \textbf{FastMAPOL version:} Prelaunch Test\\
\item
  \textbf{Data period:} March 21, 2022 (spring equinox)\\
\item
  \textbf{Reference:} Synthetic HARP2\\
\item
  \textbf{Published:} Dec 07, 2023
  (https://doi.org/10.5194/amt-16-5863-2023)
\end{itemize}

\section{Synthetic HARP2
observations}\label{synthetic-harp2-observations}

As shown in Figure~\ref{fig-validation_harp2_l1c}, A full day of
synthetic HARP2 observations is generated using a neural network forward
model for both reflectance and DoLP. Aerosol and surface properties are
prescribed from reanalysis (MERRA-2). The simulated observations provide
a realistic testbed for retrieval evaluation:

\begin{itemize}
\tightlist
\item
  Global multi-angle measurements are generated for both reflectance and
  polarization.\\
\item
  DoLP shows enhanced sensitivity to aerosol and surface properties.\\
\item
  Reflectance and DoLP provide complementary information content.
\end{itemize}

\begin{figure}

\centering{

\includegraphics[width=0.8\linewidth,height=\textheight,keepaspectratio]{../chapters/../figure/validation_harp2_l1c.png}

}

\caption{\label{fig-validation_harp2_l1c}Global HARP2 simulation (550
nm): reflectance (a--c) and DoLP (d--f) from 15 orbits on 21 March 2022,
showing three along-track viewing angles. Figure is adopted from M. Gao,
Franz, et al. (2023)}

\end{figure}%

\section{Geophysical retrievals with pixel-wise
uncertainty}\label{geophysical-retrievals-with-pixel-wise-uncertainty}

Based on the synthetic data, retrieval results as shown in
Figure~\ref{fig-validation_harp2_l2} demonstrate strong agreement with
truth and meaningful uncertainty estimates:

\begin{itemize}
\tightlist
\item
  AOD (550 nm) is retrieved with high accuracy and low spatial error.\\
\item
  Multiple geophysical variables are retrieved simultaneously, with ALH
  shown as an example.\\
\item
  Pixel-wise uncertainty reflects scene dependence, with larger values
  at low AOD and limited angular sampling.
\end{itemize}

\begin{figure}

\centering{

\includegraphics[width=0.8\linewidth,height=\textheight,keepaspectratio]{../chapters/../figure/validation_harp2_l2.png}

}

\caption{\label{fig-validation_harp2_l2}Retrieval performance: AOD (550
nm) and ALH showing retrieval, truth, and uncertainty (a--c, d--f).
Figure is adopted from M. Gao, Franz, et al. (2023)}

\end{figure}%

\section{Uncertainty closure
analysis}\label{uncertainty-closure-analysis}

Theoretical uncertainties are derived from error propagation, while
realized uncertainties are computed from retrieval--truth differences.
Their comparisons are shown in Figure~\ref{fig-validation_harp2_unc},

\begin{itemize}
\tightlist
\item
  Good agreement for aerosol optical and microphysical properties,
  especially at moderate to high AOD.\\
\item
  Decreasing uncertainty with increasing aerosol loading.\\
\item
  Underestimation of uncertainty in low-AOD conditions and for weakly
  constrained parameters.
\end{itemize}

\begin{figure}

\centering{

\includegraphics[width=0.9\linewidth,height=\textheight,keepaspectratio]{../chapters/../figure/validation_harp2_unc.png}

}

\caption{\label{fig-validation_harp2_unc}Theoretical (red) vs.~retrieved
(blue) uncertainties as a function of AOD for AOD, SSA, mr, reff, veff,
wind speed, and Chl a. AOD ranges from 0.01--0.45 (Δ=0.01). Figure is
adopted from M. Gao, Franz, et al. (2023)}

\end{figure}%

\section{Summary}\label{summary}

\begin{itemize}
\tightlist
\item
  Synthetic data provide realistic multi-angle polarimetric
  observations.\\
\item
  L2 retrievals achieve high accuracy for key variables (especially
  AOD).\\
\item
  Pixel-level uncertainties are physically meaningful and scene
  dependent.\\
\item
  Uncertainty estimates are generally consistent with actual errors,
  with known limitations.
\end{itemize}

These results define a pre-launch benchmark for HARP2 L2 validation
using DITL data.

\cleardoublepage
\phantomsection
\addcontentsline{toc}{part}{Appendices}
\appendix

\chapter{Technical Details}\label{technical-details}

This section provides detailed mathematical formulations, derivations,
and supporting information for the FastMAPOL algorithm. The material
presented here complements the main chapters and is intended for readers
who require a deeper understanding of the underlying theory,
assumptions, and implementation details.

The appendices include: - Details in radiative transfer model - Details
in the NN and automatic differentation - Detailed formulations of
aerosol size distributions and their moments\\
- Surface reflectance models for ocean and land - Multi-angle
atmospheric correction

These technical details are not required for general use or
interpretation of the FastMAPOL products but are essential for algorithm
development, validation, and advanced scientific analysis.

All equations and formulations follow the conventions introduced in the
main text unless otherwise noted.

\chapter{Radiative transfer model for PACE (PACE
Simulator)}\label{radiative-transfer-model-for-pace-pace-simulator}

The Radiative Transfer model is based on the Successive Order of
Scattering (RTSOS) method Zhai et al. (2010).

Key features: - RTSOS is capable of solving polarized radiative transfer
for both atmosphere-ocean and atmosphere-land systems. - In
atmosphere-ocean systems, it can also simulate inelastic scattering
processes, including Raman scattering and fluorescence by chlorophyll
and colored dissolved organic matter Zhai et al. (2017). - An improved
pseudos-spherical shell (IPSS) algorithm is used to improve the
simulation of radiance field over polar regions (Zhai and Hu 2022).

For this work, RTSOS is configured to simulate polarized radiance fields
for randomly generated Earth systems for the PACE instruments (Zhai and
Hu 2022). The data are then used to train neural network forward models,
which emulates the behavior of the radiative transfer system.

\textbf{Code repository} The RTSOS is made public by Dr.~Pengwang Zhai
in github.

\chapter{Neural Network Forward Model and Analytical
Jacobian}\label{neural-network-forward-model-and-analytical-jacobian}

The MAP retrievals are often computationally expensive due to their high
dimensionality and iterative nature, with multiple forward model and
Jacobian calculations. The data screening approach developed here
further increases the demand for CPU computations because the retrieval
must be repeated several times. Therefore, fast forward model and
Jacobian matrix calculations are advantageous for efficient processing,
which was a motivation for the use of NN forward models for AirHARP in
M. Gao et al. (2021). In this work, we provide the definition of the NN
and the use of automatic differentiation to compute the Jacobian matrix
analytically, as opposed to numerically through finite differencing, by
exploiting the differentiable properties of the NN models. For details
on the NN training strategies, please refer to M. Gao, Franz, et al.
(2023) and M. Gao et al. (2021).

The NN forward model developed in M. Gao et al. (2021) is a feed-forward
neural network as defined in Table~\ref{tbl-nn-forward}, where
\(\mathbf{h}_0=\mathbf{x}\) is the input layer that contains all 15
forward model parameters. Two sets of weight matrices
\(\mathbf{W}_{p+1}\) and bias vectors \(\mathbf{b}_{p+1}\) have been
determined from the NN training process for reflectance and DoLP,
respectively (M. Gao et al. 2021; M. Gao, Franz, et al. 2023).

Correspondingly, \(\mathbf{y}\) is the output layer for either
reflectance or DoLP at the four AirHARP bands.

For application to multi-angle measurements, the NN needs to be called
to simulate \(\mathbf{y}\) for each set of viewing and solar geometries
for the state vector \(\mathbf{x}\). Elements of the Jacobian matrix are
defined as follows:

\begin{equation}\phantomsection\label{eq-jacobian}{
\mathbf{K}_{mij}
=
\frac{\partial \mathbf{y}_{mi}}
{\partial \mathbf{x}_{mj}}
}\end{equation}

Here, index \(m\) indicates the viewing and solar angles, index \(i\)
indicates the wavelength, and index \(j\) indicates the state parameter.
Table~\ref{tbl-nn-forward} illustrates the structure of the NN. The
nonlinear activation function \(\Phi\) is the LeakyReLU function, which
is defined as

\begin{equation}\phantomsection\label{eq-leakyrelu}{
\Phi(\mathbf{Z})_{mi}
=
\max(0,\mathbf{Z}_{mi})
+
\alpha \times \min(0,\mathbf{Z}_{mi}).
}\end{equation}

where \(\alpha=0.01\), and \(\mathbf{Z}\) is a matrix with \(m\) and
\(i\) as its indices. The derivative with respect to each element in
\(\Phi\) is defined as

\begin{equation}\phantomsection\label{eq-D}{
\mathbf{D}_{mi}
=
\frac{\Phi(\mathbf{Z})_{mi}}
{\mathbf{Z}_{mi}}
=
\begin{cases}
1, & \mathbf{Z}_{mi}>0, \\
\alpha, & \mathbf{Z}_{mi}<0.
\end{cases}
}\end{equation}

\textbf{NN Forward Model}

\begin{longtable}[]{@{}
  >{\raggedright\arraybackslash}p{(\linewidth - 2\tabcolsep) * \real{0.5000}}
  >{\raggedright\arraybackslash}p{(\linewidth - 2\tabcolsep) * \real{0.5000}}@{}}
\caption{NN forward model.}\label{tbl-nn-forward}\tabularnewline
\toprule\noalign{}
\begin{minipage}[b]{\linewidth}\raggedright
\textbf{Layers}
\end{minipage} & \begin{minipage}[b]{\linewidth}\raggedright
\textbf{NN forward model}
\end{minipage} \\
\midrule\noalign{}
\endfirsthead
\toprule\noalign{}
\begin{minipage}[b]{\linewidth}\raggedright
\textbf{Layers}
\end{minipage} & \begin{minipage}[b]{\linewidth}\raggedright
\textbf{NN forward model}
\end{minipage} \\
\midrule\noalign{}
\endhead
\bottomrule\noalign{}
\endlastfoot
Input & \(\mathbf{h}_0=\mathbf{x}\) \\
Layer 1 &
\(\mathbf{h}_1=\Phi(\mathbf{W}_1^T\mathbf{h}_0+\mathbf{b}_1)\) \\
Layer \(p+1\) &
\(\mathbf{h}_{p+1}=\Phi(\mathbf{W}_{p+1}^T\mathbf{h}_p+\mathbf{b}_{p+1})\) \\
Output &
\(\mathbf{y}=\mathbf{W}_{k+1}^T\mathbf{h}_k+\mathbf{b}_{k+1}\) \\
\end{longtable}

\textbf{AD: Forward Mode}

\begin{longtable}[]{@{}
  >{\raggedright\arraybackslash}p{(\linewidth - 2\tabcolsep) * \real{0.5000}}
  >{\raggedright\arraybackslash}p{(\linewidth - 2\tabcolsep) * \real{0.5000}}@{}}
\caption{Jacobian matrix computed using forward-mode automatic
differentiation (AD).}\label{tbl-ad-forward}\tabularnewline
\toprule\noalign{}
\begin{minipage}[b]{\linewidth}\raggedright
\textbf{Layers}
\end{minipage} & \begin{minipage}[b]{\linewidth}\raggedright
\textbf{AD: Forward mode}
\end{minipage} \\
\midrule\noalign{}
\endfirsthead
\toprule\noalign{}
\begin{minipage}[b]{\linewidth}\raggedright
\textbf{Layers}
\end{minipage} & \begin{minipage}[b]{\linewidth}\raggedright
\textbf{AD: Forward mode}
\end{minipage} \\
\midrule\noalign{}
\endhead
\bottomrule\noalign{}
\endlastfoot
Input & \(\dot{\mathbf{h}}_{0,mij}=\delta_{ij}\) \\
Layer 1 &
\(\dot{\mathbf{h}}_{1,mij}=\mathbf{D}_{1,mi}\mathbf{W}_{1,ij}^T\) \\
Layer \(p+1\) &
\(\dot{\mathbf{h}}_{p+1,mij}=\mathbf{D}_{p+1,mi}\mathbf{W}_{p+1,il}^T\dot{\mathbf{h}}_{p,mlj}\) \\
Output &
\(\mathbf{K}_{mij}=\dot{\mathbf{y}}_{mij}=\mathbf{W}_{k+1,il}^T\dot{\mathbf{h}}_{k,mlj}\) \\
\end{longtable}

\textbf{AD: Reverse Mode}

\begin{longtable}[]{@{}
  >{\raggedright\arraybackslash}p{(\linewidth - 2\tabcolsep) * \real{0.5000}}
  >{\raggedright\arraybackslash}p{(\linewidth - 2\tabcolsep) * \real{0.5000}}@{}}
\caption{Jacobian matrix computed using reverse-mode automatic
differentiation (AD).}\label{tbl-ad-reverse}\tabularnewline
\toprule\noalign{}
\begin{minipage}[b]{\linewidth}\raggedright
\textbf{Layers}
\end{minipage} & \begin{minipage}[b]{\linewidth}\raggedright
\textbf{AD: Reverse mode}
\end{minipage} \\
\midrule\noalign{}
\endfirsthead
\toprule\noalign{}
\begin{minipage}[b]{\linewidth}\raggedright
\textbf{Layers}
\end{minipage} & \begin{minipage}[b]{\linewidth}\raggedright
\textbf{AD: Reverse mode}
\end{minipage} \\
\midrule\noalign{}
\endhead
\bottomrule\noalign{}
\endlastfoot
Output & \(\bar{\mathbf{y}}_{mij}=\mathbf{W}_{k+1,ij}^T\) \\
Layer \(p+1\) &
\(\bar{\mathbf{h}}_{p,mij}=\bar{\mathbf{h}}_{p+1,mil}\mathbf{D}_{p+1,ml}\mathbf{W}_{p+1,lj}^T\) \\
Layer 1 &
\(\bar{\mathbf{h}}_{1,mij}=\bar{\mathbf{h}}_{2,mil}\mathbf{D}_{2,ml}\mathbf{W}_{2,lj}^T\) \\
Input & \(\mathbf{K}_{mij}=\bar{\mathbf{h}}_{0,mij}\) \\
\end{longtable}

\emph{For brevity, summation over index \(l\) is implied following
Einstein notation.}

The finite difference (FD) method was used to compute the Jacobian
matrix in FastMAPOL in M. Gao et al. (2021), where the NN forward model
was called twice per input parameter under the central difference
approximation of derivatives. To reduce the computational cost, the
Jacobian matrix can be derived analytically from the NN forward model
using AD based on the chain rule of differentiation (Baydin et al.
2018). Two recursive relations are obtained to compute the Jacobian
matrix, as summarized in Table~\ref{tbl-ad-forward} and
Table~\ref{tbl-ad-reverse}. The forward mode indicates the evaluation
sequence from the first layer to the last layer, whereas the reverse
mode indicates the evaluation sequence from the last layer to the first
layer.

To represent the recursive relations, we define \(\dot{\mathbf{h}}_p\)
(tangent) and \(\bar{\mathbf{h}}_p\) (adjoint) as follows:

\[
\dot{\mathbf{h}}_{p,mij}
=
\frac{\partial \mathbf{h}_{p,mi}}
{\partial \mathbf{x}_{mj}},
\]

\[
\bar{\mathbf{h}}_{p,mij}
=
\frac{\partial \mathbf{y}_{mi}}
{\partial \mathbf{h}_{p,mj}}.
\]

Note that \(\mathbf{h}\) is defined in Table~\ref{tbl-nn-forward} as the
output from each hidden layer of the NN. The Jacobian matrix can be
represented by AD with either the tangent or the adjoint form,
corresponding to the final steps in Table~\ref{tbl-ad-forward} and
Table~\ref{tbl-ad-reverse}, respectively:

\begin{equation}\phantomsection\label{eq-forward}{
\mathbf{K}_{mij}
=
\dot{\mathbf{y}}_{mij}
}\end{equation}

\begin{equation}\phantomsection\label{eq-reverse}{
\mathbf{K}_{mij}
=
\bar{\mathbf{h}}_{0,mij}
}\end{equation}

where Equation~\ref{eq-forward} and Equation~\ref{eq-reverse} are
computed from forward- and reverse-mode AD, respectively. Forward and
reverse AD produce identical results, but differ in computational
efficiency because of the different sequence of matrix operations and
the NN architecture.

For optimal efficiency, we implemented AD directly based on the
formalism summarized in Table~\ref{tbl-ad-forward} and
Table~\ref{tbl-ad-reverse} using the PyTorch library (Paszke et al.
2019). The NN forward model is computed layer by layer, with the output
from the previous layer serving as the input to the next layer, as
summarized in Table~\ref{tbl-nn-forward}. Forward-mode AD follows the
same sequence as the NN forward model, from the first layer to the last
layer (Table~\ref{tbl-ad-forward}), whereas reverse-mode AD proceeds
from the last layer backward to the first layer
(Table~\ref{tbl-ad-reverse}). Note that AD in both modes requires the
values of matrix \(\mathbf{D}\) as defined in Equation~\ref{eq-D}, which
are determined by the output of the NN forward model at each layer.

For the NN used in this study, reverse-mode AD provides the highest
computational efficiency, as investigated further in the next section.

The AD methods provide an efficient and accurate way to compute the
Jacobian matrix, enabling substantial acceleration of retrieval
algorithms such as FastMAPOL that involve a large number of state
parameters and making them more suitable for practical applications.

\chapter{Aerosol characterization : size distribution and
representation}\label{aerosol-characterization-size-distribution-and-representation}

FastMAPOL supports several aerosol representations with increasing
complexity:

\begin{itemize}
\tightlist
\item
  \textbf{Aer-1:} five predefined log-normal modes
\item
  \textbf{Aer-2:} six predefined log-normal modes
\item
  \textbf{Aer-3:} aerosol components represented by chemical composition
\end{itemize}

These aerosol models share a common mathematical framework for size
integration, mode summation, optical properties, and bulk aerosol
products. This chapter summarizes the formulation used to connect the
retrieved state variables to physically meaningful aerosol quantities,
including optical depth, effective radius, effective variance, fine-mode
fraction, and volume fraction. In this chapter, we focused on the first
two models.

\section{Overview}\label{overview-2}

The aerosol model formulation includes the following key elements:

\begin{itemize}
\tightlist
\item
  size integration over each particle size distribution
\item
  calculation of effective radius and effective variance
\item
  multi-mode considerations
\item
  derivation of aerosol products such as fine-volume fraction (FVF) and
  fine-mode fraction (FMF)
\end{itemize}

\section{Aerosol density
distributions}\label{aerosol-density-distributions}

\subsection{Conversion between volume and number density log-normal
distributions}\label{conversion-between-volume-and-number-density-log-normal-distributions}

Both the number density and the volume density can be represented by
log-normal distributions. The two forms are equivalent and can be
converted into each other by changing the mean radius and the
normalization constant.

For a monomodal system, the volume and number size distributions are
written as

\[
\frac{d V(r)}{d \ln r}= \frac{V_{0}}{\sqrt{2 \pi} \sigma_{v}} \exp\left[-\frac{(\ln r-\ln r_{v})^2}{2 \sigma_{v}^2}\right],
\]

\[
\frac{d N(r)}{d \ln r}= \frac{N_{0}}{\sqrt{2 \pi} \sigma_{n}} \exp\left[-\frac{(\ln r-\ln r_{n})^2}{2 \sigma_{n}^2}\right].
\]

Here, \(r_v\) and \(r_n\) are the mean radii for the volume and number
distributions, respectively, and \(\sigma_v\) and \(\sigma_n\) are the
corresponding logarithmic standard deviations. The formulation follows
Xu (2016, Appendix A6), with the corrected definition of \(\sigma\) as
confirmed by Pengwang Zhai through personal communication with Feng Xu.

Since particle volume and number are related by

\[
V=\frac{4\pi}{3}Nr^3,
\]

the corresponding differential distributions satisfy

\[
\frac{d V(r)}{d \ln r}=\frac{4\pi}{3} r^3 \frac{d N(r)}{d \ln r}.
\]

Using this relation, the following identities are obtained:

\[
\sigma_{n}=\sigma_{v}=\sigma,
\]

\[
r_{n}=r_{v} \exp(-3 \sigma^2),
\]

\[
N_{0}=V_{0} \frac{3}{4 \pi r_{v}^3}\exp(4.5 \sigma^2),
\]

\[
N_{0}=V_{0} \frac{3}{4 \pi r_{n}^3}\exp(-4.5 \sigma^2).
\]

These results are consistent with the derivation given by R.G. Grainger.
The last two expressions provide an internal consistency check.

\section{Optical depth}\label{optical-depth}

Using the conversion above, the column-integrated number density can be
expressed in terms of the column-integrated volume density as

\[
N_{0}=V_{0} \frac{3}{4 \pi r_{v}^3}\exp(4.5 \sigma^2).
\]

The aerosol optical depth is defined as

\[
\tau_a=\sum_{i=1}^{N_c} C_{ext,i} N_i,
\]

where \(C_{ext,i}\) is the extinction cross section of component \(i\)
in units of \(\mu m^2\), and \(N_i\) is the corresponding column number
density in units of \(\mu m^{-2}\).

Here, column number density is equivalent to vertically integrated
number density. If the local number density is \(n\) in units of
\(\mu m^{-3}\) and the layer thickness is \(z\) in units of \(\mu m\),
then

\[
N = n z.
\]

Using column volume density, the optical depth can be rewritten as

\[
\tau_a = \sum_{i=1}^{N_c} C_{ext,i} N_i,
\]

\[
\tau_a = \sum_{i=1}^{N_c} C_{ext,i} V_i \frac{3}{4 \pi r_{v}^3}\exp(4.5 \sigma^2),
\]

\[
\tau_a = \sum_{i=1}^{N_c} \left[C_{ext,i} \frac{3}{4 \pi r_{v}^3}\exp(4.5 \sigma^2)\right] V_i.
\]

The unit of column volume density is \(\mu m^3 \mu m^{-2}\), which
reduces to \(\mu m\).

If a maximum optical depth \(\tau_{max}\) is prescribed, then the
corresponding maximum allowed volume density for each component is

\[
V_i=\tau_{max}\left[C_{ext,i} \frac{3}{4 \pi r_{v}^3}\exp(4.5 \sigma^2)\right]^{-1},
\]

or equivalently

\[
V_i=\tau_{max} \exp(4.5 \sigma^2) \frac{4 \pi r_{v}^3}{3 C_{ext,i}}.
\]

A practical bound on \(V_i\) can be estimated at a reference wavelength
such as \(550 \, nm\) by computing extinction cross sections over the
expected refractive-index range for each mode. For the lower bound, one
may use a real refractive index close to that of water at \(550 \, nm\),
approximately 1.3.

In general, increasing the imaginary part of the refractive index
increases the extinction cross section. For coarse particles, however,
the extinction is often much less sensitive to refractive index than to
particle size.

\section{Convention for size
distributions}\label{convention-for-size-distributions}

The size distributions may be defined with respect to either \(r\) or
\(\ln r\). These forms are related by

\[
\frac{d N(r)}{d r}=\frac{1}{r}\frac{d N(r)}{d \ln r},
\]

\[
\frac{d V(r)}{d r}=\frac{1}{r}\frac{d V(r)}{d \ln r}.
\]

This distinction is important when comparing formulations from different
references or numerical implementations.

\section{Moments of the particle size
distribution}\label{moments-of-the-particle-size-distribution}

For later use, define the \(k\)th moment of the number size distribution
as

\[
m_k=\int_{0}^{\infty} r^k \frac{d N(r)}{d r} \, d r,
\]

which is equivalently

\[
m_k=\int_{-\infty}^{\infty} r^k \frac{d N(r)}{d \ln r} \, d \ln r.
\]

For a log-normal number distribution,

\[
m_k=N_0 r_n^k \exp\left(\frac{k^2}{2} \sigma^2\right).
\]

The dimension is consistent: \(m_k\) has units of \(r^k\), while
\(\sigma\) is dimensionless.

Substituting the number-volume conversion into the moment expression
gives

\[
m_k=N_0 r_n^k \exp\left(\frac{k^2}{2} \sigma^2\right),
\]

\[
m_k=V_{0} \frac{3}{4 \pi r_{v}^3}\exp(4.5 \sigma^2)r^k_{v} \exp(-3k \sigma^2) \exp\left(\frac{k^2}{2} \sigma^2\right),
\]

and therefore

\[
m_k=V_{0} \frac{3}{4 \pi} r^{k-3}_v\exp\left[\frac{(k-3)^2\sigma^2}{2} \right].
\]

This result preserves dimensional consistency because \(V_0\) carries
the dimension of \(r^3\).

A useful special case is \(k=3\), for which

\[
m_3=\frac{3V_0}{4\pi}.
\]

Thus, the third moment depends directly on the total particle volume.

\section{\texorpdfstring{Effective radius and effective variance in
terms of
\(dN/d\ln r\)}{Effective radius and effective variance in terms of dN/d\textbackslash ln r}}\label{effective-radius-and-effective-variance-in-terms-of-dndln-r}

The effective radius, or area-weighted mean radius, is defined as

\[
r_{eff}=\frac{\int_{r_1}^{r_2} r \, r^2 (dN(r)/dr) \, dr}{\int_{r_1}^{r_2} r^2 (dN(r)/dr) \, dr}
=\frac{m_3}{m_2}.
\]

The weighting by particle area reflects the fact that radiative
extinction is approximately proportional to particle cross-sectional
area.

The effective variance is defined as

\[
\nu_{eff}=\frac{\int_{r_1}^{r_2}(r-r_{eff})^2 r^2 (dN(r)/dr) \, dr}
{r^2_{eff} \int_{r_1}^{r_2}r^2 (dN(r)/dr) \, dr},
\]

which can be written in terms of moments as

\[
\nu_{eff}= \frac{m_4 m_2}{m_3^2}-1.
\]

This expression is equally valid for distributions written in terms of
\(dN(r)/d\ln r\).

For a log-normal number distribution

\[
\frac{d N(r)}{d \ln r}=\frac{N_0}{\sqrt{2 \pi} \sigma} \exp\left[-\frac{(\ln r-\ln r_n)^2}{2 \sigma^2}\right],
\]

the effective radius and effective variance are

\[
r_{eff}=r_n \exp\left(\frac{5}{2} \sigma^2\right),
\]

\[
\nu_{eff}=\exp (\sigma^2)-1.
\]

Using the relation between number-mean and volume-mean radius,

\[
r_{eff}=r_n \exp(2.5 \sigma^2),
\]

\[
r_{eff}=r_v \exp(-3 \sigma^2) \exp(2.5 \sigma^2),
\]

\[
r_{eff}=r_v \exp(-0.5 \sigma^2).
\]

Thus, the effective radius is always larger than the number-mean radius
and smaller than the volume-mean radius.

For small \(\sigma^2\), the approximations

\[
\nu_{eff}=\exp (\sigma^2)-1 \approx \sigma^2
\]

and

\[
r_{eff}= r_v \exp(-0.5 \sigma^2) \approx r_v
\]

show that the effective variance is approximately the square of the
logarithmic standard deviation, while the effective radius approaches
the volume-mean radius.

\section{\texorpdfstring{Effective radius and variance in terms of
\(dV/d\ln r\)}{Effective radius and variance in terms of dV/d\textbackslash ln r}}\label{effective-radius-and-variance-in-terms-of-dvdln-r}

Because the retrieval state is often expressed in terms of volume
density, it is useful to write effective radius and effective variance
directly in terms of the volume distribution.

Starting from

\[
\frac{d V(r)}{d \ln r}=\frac{4\pi}{3} r^3 \frac{d N(r)}{d \ln r},
\]

the effective radius becomes

\[
r_{eff}=\frac{\int_{\ln r_1}^{\ln r_2} (dV(r)/d\ln r) \, d\ln r}
{\int_{\ln r_1}^{\ln r_2} r^{-1} (dV(r)/d\ln r) \, d\ln r}.
\]

Define the moments of the volume distribution as

\[
m'_k = \int_{\ln r_1}^{\ln r_2} r^k \left(\frac{dV(r)}{d\ln r}\right) d\ln r.
\]

Then

\[
r_{eff}= \frac{m'_0}{m'_{-1}}.
\]

Similarly, the effective variance becomes

\[
\nu_{eff}=\frac{\int_{\ln r_1}^{\ln r_2}(r-r_{eff})^2 r^{-1} (dV(r)/d\ln r) \, d\ln r}
{r^2_{eff} \int_{\ln r_1}^{\ln r_2} r^{-1} (dV(r)/d\ln r) \, d\ln r},
\]

which simplifies to

\[
\nu_{eff}= \frac{m'_{-1} m'_1}{m'^2_0}-1.
\]

For discrete numerical implementation, the moment may be written
schematically as

\[
m'_k \approx \sum_i r(i)^k v(i,j),
\]

where \(i\) indexes the size grid and \(j\) indexes pixel or retrieval
location. The constant increment \(d\ln r\) is omitted here for brevity.

For a log-normal volume distribution,

\[
m'_k =V_{0} \frac{3}{4 \pi} r^{k}_v\exp\left[\frac{k^2\sigma^2}{2} \right].
\]

For a multimode distribution, the effective quantities become

\[
r_{eff}=\frac{\sum_i m_{3,i}}{\sum_i m_{2,i}}
\]

and

\[
\nu_{eff}=\frac{\sum_i m_{4,i} \sum_i m_{2,i}}{(\sum_i m_{3,i})^2}-1
\]

when expressed through number-density moments, or equivalently

\[
r_{eff}=\frac{\sum_i m'_{0,i}}{\sum_i m'_{-1,i}},
\]

\[
\nu_{eff}=\frac{\sum_i m'_{1,i} \sum_i m'_{-1,i}}{(\sum_i m'_{0,i})^2}-1
\]

when expressed through volume-density moments.

\section{Multi-mode distributions}\label{multi-mode-distributions}

When the aerosol is represented as a sum of log-normal modes,

\[
\frac{d V(r)}{d \ln r}= \sum_i \frac{V_{0,i}}{\sqrt{2 \pi} \sigma_{v,i}} \exp\left[-\frac{(\ln r-\ln r_{v,i})^2}{2 \sigma_{v,i}^2}\right],
\]

\[
\frac{d N(r)}{d \ln r}= \sum_i \frac{N_{0,i}}{\sqrt{2 \pi} \sigma_{n,i}} \exp\left[-\frac{(\ln r-\ln r_{n,i})^2}{2 \sigma_{n,i}^2}\right].
\]

For each mode, the conversion relations are

\[
\sigma_{n,i}=\sigma_{v,i}=\sigma_i,
\]

\[
r_{n,i}=r_{v,i} \exp(-3 \sigma_{i}^2),
\]

\[
N_{0,i}=V_{0,i} \frac{3}{4 \pi r_{v,i}^3}\exp(4.5 \sigma_{i}^2),
\]

\[
N_{0,i}=V_{0,i} \frac{3}{4 \pi r_{n,i}^3}\exp(-4.5 \sigma_{i}^2).
\]

The distinction between number weighting and volume weighting remains
important, since the same physical aerosol can have different mean radii
depending on which representation is used.

For the volume-mode parameterization adopted from O. Dubovik et al.
(2011) and Xu et al. (2016), the nominal mode radii are

\[
0.1,\ 0.1732,\ 0.3,\ 1.0,\ 2.9 \ \mu m
\]

with logarithmic standard deviations

\[
0.35,\ 0.35,\ 0.35,\ 0.5,\ 0.5.
\]

The impact of the number of modes are further discussed by Fu:2020aa. We
adopted this approach as default aerosol model (Aer-1). Meng Gao et al.
(2018) add an additional coarse mode to deal with the larger aerosol
size when using the SWIR bands from RSP. The nomial radii and standard
deviations are 8.4 \(\mu m\) and 0.5, which is called (Aer-2).

\section{Averaged extinction cross
section}\label{averaged-extinction-cross-section}

When extinction is averaged over a volume-based aerosol model, number
density must first be recovered from the volume density.

For the same total particle volume, a fine mode contains far more
particles than a coarse mode. As a result, fine-mode particles can
dominate number-weighted optical properties even when their volume
contribution is modest.

For two modes, the total extinction can be written as

\[
C_{ext,t}=C_{ext,f} n_f+C_{ext,c} n_c,
\]

or in terms of volume density,

\[
C_{ext,t}= C_{ext,f} \frac{3 V_f}{4\pi r^3_{v,f}}\exp(4.5 \sigma^2_f)+
C_{ext,c} \frac{3 V_c}{4\pi r^3_{v,c}}\exp(4.5 \sigma^2_c).
\]

To compare fine- and coarse-mode contributions, consider the ratio

\[
\frac{C_{ext,f} \frac{3 V_f}{4\pi r^3_{v,f}}\exp(4.5 \sigma^2_f)}
{C_{ext,c} \frac{3 V_c}{4\pi r^3_{v,c}}\exp(4.5 \sigma^2_c)}.
\]

If the logarithmic standard deviations are assumed equal, this ratio
simplifies to an expression proportional to

\[
\frac{V_f}{V_c}\frac{Q_{ext,f} r_{v,c}}{Q_{ext,c} r_{v,f}},
\]

showing explicitly how fine-mode dominance arises from both larger
number concentration per unit volume and the size dependence of
extinction efficiency.

\section{Effective radius and variance for multimode
aerosols}\label{effective-radius-and-variance-for-multimode-aerosols}

The total number density and total volume density are

\[
N_0=\sum_i N_{0,i},
\]

\[
V_0=\sum_i V_{0,i}.
\]

The effective radius and effective variance of the combined distribution
are

\[
r_{eff}=\frac{\sum_i m_{3,i}}{\sum_i m_{2,i}},
\]

\[
\nu_{eff}=\frac{\sum_i m_{4,i} \sum_i m_{2,i}}{(\sum_i m_{3,i})^2}-1.
\]

\subsection{Sensitivity to component
volume}\label{sensitivity-to-component-volume}

To propagate retrieval uncertainty from mode volumes to bulk effective
properties, derivatives with respect to \(V_j\) are required.

For effective radius,

\[
\frac{\partial r_{eff}}{\partial V_{j}} =
\frac{1}{\sum_i  m_{2,i}}\frac{\partial m_{3,j}}{\partial V_j}- 
\frac{\sum_i m_{3,i}}{(\sum_i m_{2,i})^2}\frac{\partial m_{2,j}}{\partial V_j},
\]

which can be written as

\[
\frac{\partial r_{eff}}{\partial V_{j}} =
r_{eff} \left(\frac{m'_{3,j}}{\sum_i  m_{3,i}}-\frac{m'_{2,j}}{\sum_i  m_{2,i}}\right).
\]

For effective variance,

\[
\frac{\partial \nu_{eff}}{\partial V_j}=
\frac{\sum_i m_{2,i}}{(\sum_i m_{3,i})^2}\frac{\partial m_{4,j}}{\partial V_j}+
\frac{\sum_i m_{4,i}}{(\sum_i m_{3,i})^2}\frac{\partial m_{2,j}}{\partial V_j}
-2\frac{\sum_i m_{4,i} \sum_i m_{2,i}}{(\sum_i m_{3,i})^3}\frac{\partial m_{3,j}}{\partial V_j},
\]

or equivalently

\[
\frac{\partial \nu_{eff}}{\partial V_j}=
(\nu_{eff}+1) \left(\frac{m'_{4,j}}{\sum_i  m_{4,i}}+\frac{m'_{2,j}}{\sum_i m_{2,i}} 
-2\frac{m'_{3,j}}{\sum_i  m_{3,i}}\right).
\]

Here the derivative of the moment with respect to component volume is

\[
m'_{k,j}=\frac{\partial m_{k,j}}{\partial V_j}=\frac{m_{k,j}}{V_j},
\]

which gives

\[
m'_{k,j}=\frac{3}{4 \pi}
r^{k-3}_{v,j} \exp\left[\frac{(k-3)^2}{2} \sigma_j^2\right].
\]

For \(k=3\), this derivative is constant because \(m_3\) depends
linearly on volume and not on radius.

\section{Physical interpretation of low-order
moments}\label{physical-interpretation-of-low-order-moments}

For a volume-based log-normal mode, the low-order moments are

\[
m_{2,i}=\frac{3}{4\pi} V_i r_{v,i}^{-1} \exp(\sigma_i^2/2),
\]

\[
m_{3,i}=\frac{3}{4\pi} V_i,
\]

\[
m_{4,i}=\frac{3}{4\pi} V_i r_{v,i} \exp(\sigma_i^2/2).
\]

For a single mode, these lead directly to

\[
r_{eff}=\frac{m_3}{m_2}=r_v \exp(-\sigma^2/2),
\]

\[
\nu_{eff}=\frac{m_4 m_2}{m_3^2}-1=\exp (\sigma^2)-1.
\]

For a multimode aerosol,

\[
\sum_i m_{3,i}=\frac{3}{4\pi}\sum_i V_i=\frac{3}{4\pi} V_0,
\]

so the summed third moment is proportional to the total particle volume.

Within a fine-mode or coarse-mode subset, \(m_2\) emphasizes smaller
particles, whereas \(m_4\) emphasizes larger particles. This makes
effective radius and effective variance sensitive to the relative
partitioning among submodes, even when the logarithmic standard
deviation is fixed within that subset.

\section{Fine and coarse mode
partitions}\label{fine-and-coarse-mode-partitions}

To discuss fine-mode and coarse-mode aerosol separately, the mode
summation may be restricted to selected ranges.

For the Aer-1 and Aer-2 parameterization adopted here:

\begin{itemize}
\tightlist
\item
  \textbf{fine mode:} modes 1 to 3, with volume-mean radii from \(0.1\)
  to \(0.3 \, \mu m\)
\item
  \textbf{coarse mode:} modes 4 and above, with volume-mean radii from
  \(1.0\), \(2.9, (and 8.4) \mu m\)
\end{itemize}

Although the first three fine modes share the same logarithmic standard
deviation, their effective variance can still change because the bulk
quantity depends on the relative partitioning among the mode volumes.

\section{Phase-matrix averaging}\label{phase-matrix-averaging}

The total number density may be represented as a weighted sum of
component size distributions,

\[
n(r)=\sum_i c_i n_i(r),
\]

with

\[
\sum_i c_i =1.
\]

The corresponding size-averaged phase matrix is then

\[
P(\Omega)_{avg}=\int d r \, n(r) P(\Omega;r),
\]

which can be expanded as

\[
P(\Omega)_{avg}= \sum_i c_i \int d r \, n_i(r) P(\Omega;r)
\]

and written as

\[
P(\Omega)_{avg}= \sum_i c_i P_i (\Omega).
\]

If \(P_i\) denotes the size-averaged phase matrix of component \(i\),
then the total phase matrix is a weighted sum over components. If
normalized phase matrices \(\tilde{P}_i\) are used, then the weighting
must be consistent with the corresponding scattering cross sections.

In practice, this requires confirming the normalization convention of
the scattering subroutine, such as \texttt{spher\_interface}.

\section{Additional notes on units}\label{additional-notes-on-units}

The column number density may also be written as

\[
N_0=n_0 \delta l,
\]

where \(\delta l\) is the aerosol layer thickness.

In log-normal expressions, the absolute unit of radius does not matter
as long as \(r\) and the mean radius use the same unit, since

\[
\ln r-\ln r_m=\ln(r/r_m).
\]

Both the logarithmic variance and the effective variance are
dimensionless.

\section{Optical depth at different
wavelengths}\label{optical-depth-at-different-wavelengths}

If the aerosol loading is parameterized at a reference wavelength, such
as \(555 \, nm\), then the corresponding number density can be defined
as

\[
n_i=\frac{\tau_{555nm}}{C_{ext,555nm}}.
\]

The optical depth at any other wavelength is then

\[
\tau_{\lambda}=n_i C_{ext,\lambda}.
\]

This relation is useful when the state variable is defined through
reference-wavelength optical depth while the forward model requires
spectrally varying extinction.

\section{Aerosol products}\label{aerosol-products}

The aerosol model supports the derivation of bulk products from the
retrieved mode volumes and optical properties. These include, for
example:

\begin{itemize}
\tightlist
\item
  \textbf{FMF:} fine-mode fraction, typically defined from fine-mode
  optical depth relative to total optical depth at a reference
  wavelength
\item
  \textbf{FVF:} fine-volume fraction, defined from fine-mode particle
  volume relative to total particle volume
\item
  effective radius and effective variance for fine, coarse, or total
  aerosol
\item
  spectrally resolved optical depth and related intensive properties
\end{itemize}

The exact definition of each product depends on the chosen aerosol
representation and on whether the partitioning is performed by mode, by
size, or by composition.

\section{Summary}\label{summary-1}

The aerosol model is built on a consistent log-normal framework that
links retrieved mode volumes to number density, optical depth, size
moments, and bulk aerosol properties. The volume-based representation is
particularly convenient for retrieval, while the number-based
representation remains useful for interpreting extinction and other
radiatively weighted quantities. The relationships summarized here
provide the basis for size integration, mode summation, and the
calculation of operational aerosol products in FastMAPOL.

\chapter{Uncertainty estimation and angular
correlations}\label{uncertainty-estimation-and-angular-correlations}

In this section, we discuss the general retrieval approach considering
correlation between angles following (M. Gao, Knobelspiesse, et al.
2023). The same maximum-likelihood approach is used to retrieve the
state parameters in FastMAPOL by minimizing a cost function that
represents the difference between the measurements and forward-model
calculations (Rodgers 2000) but in a matrix form:

\begin{equation}\phantomsection\label{eq-chi2}{
\chi^2
=
\frac{1}{N}
\mathbf{y}^{T}
\mathbf{S}_{\epsilon}^{-1}
\mathbf{y},
}\end{equation}

where

\[
\mathbf{y}
=
\mathbf{m}
-
\mathbf{f}(\mathbf{x})
\]

is the residual vector between the measurement vector \(\mathbf{m}\) and
forward model \(\mathbf{f}\) for the state vector \(\mathbf{x}\). The
measurement vector \(\mathbf{m}\) includes both reflectance (\(\rho_t\))
and DoLP (\(P_t\)), where the subscript \(t\) indicates the total signal
measured by the instrument. The total number of measurements, \(N\), at
each pixel includes contributions from both reflectance and DoLP (M. Gao
et al. 2021; Meng Gao et al. 2021).

The error covariance matrix \(\mathbf{S}_{\epsilon}\) in
Equation~\ref{eq-chi2} specifies both the uncertainty of each
measurement and correlations among measurements at the same pixel. It is
a symmetric matrix defined as

\begin{equation}\phantomsection\label{eq-covariance-definition}{
S_{\epsilon;i,j}
=
\mathbb{E}
\left[
(y_i-\mathbb{E}[y_i])
(y_j-\mathbb{E}[y_j])
\right],
}\end{equation}

where \(i\) and \(j\) indicate measurements at different viewing angles
and spectral bands, and \(\mathbb{E}\) denotes the expectation value.

To represent angular uncertainty correlations, an autoregressive model
of order 1, denoted AR(1), is used. This approach was previously applied
to RSP data (Knobelspiesse et al. 2012) and is adopted here for HARP
data. AR(1) represents a linear Markov process with the error covariance
matrix specified as

\begin{equation}\phantomsection\label{eq-covariance}{
S_{\epsilon;i,j}
=
\begin{cases}
\sigma_{t,i}^2,
& i=j, \\[4pt]
\sigma_{c,i}\sigma_{c,j}
r^{\Delta_\theta |i-j|},
& i\ne j \text{ and at the same band and polarization state}, \\[4pt]
0,
& \text{otherwise}.
\end{cases}
}\end{equation}

where \(\sigma_t\) is the total uncertainty, including both random noise
and calibration uncertainty (\(\sigma_c\)). Only \(\sigma_c\) is assumed
to be correlated among measurements at different viewing angles.

The relative contributions of random and calibration uncertainties may
differ between reflectance and polarized signals (Knobelspiesse et al.
2019). For synthetic data generated directly using the forward model,
the contribution of forward-model uncertainty is not considered.
\(\Delta_\theta\) is the average angular grid spacing and depends on the
spectral channel. We model the correlation properties using
\(\Delta_\theta\) estimated from the viewing angles in the along-track
direction to better represent the stripe-filter characteristics used for
HARP angular measurements.

The average \(\Delta_\theta\) is approximately \(6.0^\circ\) for AirHARP
and \(12^\circ\) for HARP2 at 440, 550, and 870 nm, and approximately
\(2.0^\circ\) at 670 nm for both HARP instruments. The parameter \(r\)
in Equation~\ref{eq-covariance} is the correlation parameter, with a
value between 0 and 1. For uncertainties with more complex structures, a
general autoregressive moving-average (ARMA) model can be used
(Priestley 1983). However, analysis of retrieval results from real
AirHARP measurements indicates that the AR(1) model provides a useful
representation for most cases.

To provide a more intuitive representation of strong correlations as
\(r\) approaches 1, we define a correlation angle \(\theta_c\) according
to

\begin{equation}\phantomsection\label{eq-correlation-angle}{
r^{\Delta_\theta |i-j|}
=
e^{-\Delta_\theta |i-j|/\theta_c}.
}\end{equation}

Therefore, \(\theta_c\) represents the angular separation over which the
magnitude of the correlation is reduced by a factor of \(e\). The
correlation angle can be derived from \(r\) as

\begin{equation}\phantomsection\label{eq-thetac2r}{
\theta_c
=
-\frac{1}{\ln r}.
}\end{equation}

\section{Uncertainty Quantifization}\label{uncertainty-quantifization}

Pixel-wise retrieval uncertainty can be quantified by propagating
measurement and forward-model uncertainties into the retrieval parameter
space (Rodgers 2000):

\begin{equation}\phantomsection\label{eq-S}{
\mathbf{S}^{-1}
=
\mathbf{K}^{T}
\mathbf{S}_{\epsilon}^{-1}
\mathbf{K}
+
\mathbf{S}_{a}^{-1},
}\end{equation}

where the Jacobian matrix \(\mathbf{K}\) contains the partial
derivatives of the modeled measurements with respect to the retrieval
parameters.

Each retrieval parameter is constrained to a finite permitted range,
which imposes an implicit a priori constraint on the retrieval. To
represent its influence on retrieval uncertainties, the a priori error
covariance matrix \(\mathbf{S}_a\) in Equation~\ref{eq-S} is assumed to
be diagonal, with the a priori uncertainty for each state parameter
approximated by its permitted retrieval range (M. Gao et al. 2022).

The uncertainties are defined as the standard deviation (\(1\sigma\))
around the retrieval solution and are estimated from the square roots of
the diagonal elements of \(\mathbf{S}\). Uncertainties for variables
that are functions of the retrieved state parameters can similarly be
derived from \(\mathbf{S}\) and the corresponding derivatives.

Because FastMAPOL retrieves a relatively large number of state
parameters, direct evaluation of retrieval uncertainties can be
computationally expensive. The uncertainty calculation is accelerated
using automatic differentiation applied to the neural network forward
models (M. Gao et al. 2022). For example, the uncertainty in remote
sensing reflectance (\(R_{rs}\)) can be propagated using automatic
differentiation through the neural networks used for atmospheric and
BRDF corrections.

The retrieval uncertainties estimated through error propagation in
Equation~\ref{eq-S}, hereafter referred to as \textbf{theoretical
retrieval uncertainties}, represent an optimalized local estimate and
rely on assumptions such as successful convergence of the retrieval
parameters to the appropriate solution (Sayer et al. 2020; M. Gao et al.
2022). If the retrieval solution differs substantially from the true
state, both the retrieved state and its associated Jacobian may be less
representative of the true conditions, potentially leading to inaccurate
uncertainty estimates.

For synthetic retrieval experiments, the retrieval error can be
calculated directly as the difference between the retrieved and true
values. The statistical distribution of these retrieval errors is
referred to here as the \textbf{real retrieval uncertainty}. Comparing
theoretical and real retrieval uncertainties provides a means of
evaluating the optimal and actual performance of the retrieval
algorithm.

The Monte Carlo error propagation (MCEP) method is used for this
comparison (M. Gao et al. 2022). MCEP samples retrieval errors from the
theoretical retrieval uncertainty distributions and directly compares
these sampled error distributions with errors obtained from the actual
retrievals. Multiple sets of random samples can be generated to evaluate
the influence of sample size on the estimated uncertainty statistics (M.
Gao et al. 2022).

\section{Eigenvector decomposition of the error covariance
matrix}\label{sec-transform}

An error covariance matrix with non-diagonal terms can be challenging to
implement efficiently in optimization algorithms, which often operate
most conveniently in a space with uncorrelated measurement errors.
Correlated measurement uncertainties can also be difficult to interpret
because individual uncertainty components are associated with multiple
measurements. To address these issues, the measurements can be
transformed into a new space in which the error covariance matrix is
diagonal.

Eigenvector decomposition is applied to the error covariance matrix
(Rodgers 2000):

\begin{equation}\phantomsection\label{eq-svd}{
\mathbf{S}_{\epsilon}
=
\mathbf{U}^{T}
\mathbf{D}_{\epsilon}
\mathbf{U},
}\end{equation}

where \(\mathbf{D}_{\epsilon}\) is a positive diagonal matrix containing
the eigenvalues of \(\mathbf{S}_{\epsilon}\) and \(\mathbf{U}\) is a
unitary matrix.

Based on Equation~\ref{eq-svd}, the cost function in
Equation~\ref{eq-chi2} and the error propagation in Equation~\ref{eq-S}
can be written as

\begin{equation}\phantomsection\label{eq-chi2-new}{
\chi^2
=
\frac{1}{N}
\mathbf{y}'^{T}
\mathbf{D}_{\epsilon}^{-1}
\mathbf{y}',
}\end{equation}

and

\begin{equation}\phantomsection\label{eq-S-new}{
\mathbf{S}^{-1}
=
\mathbf{K}'^{T}
\mathbf{D}_{\epsilon}^{-1}
\mathbf{K}'
+
\mathbf{S}_{a}^{-1},
}\end{equation}

where the original residual vector \(\mathbf{y}\) is transformed to a
new residual vector \(\mathbf{y}'\) with uncorrelated uncertainties,
together with the corresponding Jacobian matrix:

\begin{equation}\phantomsection\label{eq-y-transform}{
\mathbf{y}'
=
\mathbf{U}\mathbf{y},
}\end{equation}

\begin{equation}\phantomsection\label{eq-K-transform}{
\mathbf{K}'
=
\mathbf{U}\mathbf{K}.
}\end{equation}

The transformations in Equation~\ref{eq-chi2-new} and
Equation~\ref{eq-S-new} allow the retrieval to operate with diagonal
uncertainty matrices and provide several advantages:

\begin{enumerate}
\def\labelenumi{\arabic{enumi}.}
\item
  \textbf{Clear representation of measurement uncertainty.}\\
  The diagonal elements of \(\mathbf{D}_{\epsilon}\) represent the
  uncertainties of the transformed measurements \(\mathbf{y}'\),
  providing insight into how measurement accuracy is affected by
  correlations.
\item
  \textbf{Efficient minimization of the retrieval cost function.}\\
  Equation~\ref{eq-chi2-new} represents the cost function for the
  uncorrelated transformed measurements \(\mathbf{y}'\) and can
  therefore be used directly with conventional optimization algorithms,
  such as the subspace trust-region interior-reflective (STIR) algorithm
  (Branch, Coleman, and Li 1999), as implemented in FastMAPOL (M. Gao et
  al. 2021, 2022; Meng Gao et al. 2021).
\item
  \textbf{Generation of correlated errors.}\\
  To investigate and visualize angular uncertainty correlations, errors
  can first be generated in the transformed \(\mathbf{y}'\) space by
  drawing random samples from normal distributions whose variances are
  defined by the eigenvalues in \(\mathbf{D}_{\epsilon}\). These errors
  are then transformed back into the original measurement space
  according to

  \begin{equation}\phantomsection\label{eq-y-inverse-transform}{
  \mathbf{y}
  =
  \mathbf{U}^{T}\mathbf{y}'.
  }\end{equation}

  This procedure enables synthetic measurement errors with prescribed
  angular correlation structures to be generated and added to simulated
  measurements.
\end{enumerate}

\section{Correlation strength estimation using
autocorrelation}\label{sec-ac}

Autocorrelation provides a useful measure of correlation in a discrete
data sequence and is defined as (Priestley 1983)

\begin{equation}\phantomsection\label{eq-autocorrelation}{
R_{i,j}
=
\mathbb{E}[y_i y_j],
}\end{equation}

where \(i\) and \(j\) are two indices of the dataset. Comparing
Equation~\ref{eq-autocorrelation} with
Equation~\ref{eq-covariance-definition}, the autocorrelation is
equivalent to the autocovariance when \(\mathbb{E}[y_i]=0\).

This method can be applied to simulated noise and to retrieval fitting
residuals. However, the mean and variance of the fitting residuals can
vary with viewing angle. Such signals are non-stationary and are
difficult to represent directly using AR models (Priestley 1983). To
reduce this effect, the original residual vector \(\mathbf{y}\) is
processed by removing its mean and normalizing by its standard
deviation. The resulting normalized residual is denoted by
\(\tilde{\mathbf{y}}\).

For measurements within the same spectral band and polarization state,
the autocorrelation function of the normalized residuals is equivalent
to the covariance defined in Equation~\ref{eq-covariance}:

\begin{equation}\phantomsection\label{eq-rk}{
\tilde{R}_k
=
\mathbb{E}
\left[
\tilde{y}_i
\tilde{y}_{i+k}
\right]
=
r^{\Delta_\theta k}.
}\end{equation}

The correlation strength can therefore be estimated by analyzing the
residuals between the measurements and forward-model calculations. The
autocorrelation function is averaged over multiple pixels to reduce
statistical uncertainty. The correlation parameter can then be estimated
as

\begin{equation}\phantomsection\label{eq-r-from-ac}{
r
=
\left(\tilde{R}_1\right)^{1/\Delta_\theta}.
}\end{equation}

The corresponding correlation angle \(\theta_c\) is calculated from
\(r\) using Equation~\ref{eq-thetac2r}.

A partial autocorrelation function can also be calculated from the data
sequence. Partial autocorrelation removes correlations associated with
intermediate lags and can therefore be used to evaluate whether an AR(1)
model adequately represents the uncertainty structure (Priestley 1983).
If AR(1) is sufficient, only the first-order lag should remain prominent
beyond the zero-order term. The Python packages StatsModels (Seabold and
Perktold 2010) and SciPy (Virtanen et al. 2018) are used to conduct the
autocorrelation analyses.

\chapter{Ocean Bio-Optical Model}\label{ocean-bio-optical-model}

\section{Overview}\label{overview-3}

FastMAPOL retrieves ocean optical properties by coupling the aerosol
retrieval with a forward radiative transfer model of ocean inherent
optical properties (IOPs). The ocean optical properties are
parameterized using bio-optical models that describe the spectral
absorption and scattering of seawater constituents.

Three parameterizations are considered:

\begin{longtable}[]{@{}lll@{}}
\toprule\noalign{}
Model & Parameters & Application \\
\midrule\noalign{}
\endhead
\bottomrule\noalign{}
\endlastfoot
Bio-1 & 1 parameter & Open ocean (Case-1 waters) \\
Bio-2 & 5 parameters & Coastal and optically complex waters \\
Bio-3 & 3 parameters & Optimized coastal parameterization \\
\end{longtable}

These models describe the spectral absorption and scattering properties
of four main water constituents:

\begin{itemize}
\tightlist
\item
  pure seawater
\item
  phytoplankton
\item
  non-algal particles (NAP)
\item
  colored dissolved organic matter (CDOM)
\end{itemize}

The total inherent optical properties are constructed from the
contributions of these components.

\section{Bio-1 Optical Model (Open
Ocean)}\label{bio-1-optical-model-open-ocean}

The \textbf{Bio-1 model} is a simplified bio-optical parameterization
designed for \textbf{open ocean waters (Case-1 waters)}. The model is
derived from the generalized Bio-2 model by imposing constraints on its
parameters (Zhai et al. (2015), Zhai et al. (2017)).

In this model, the optical properties are parameterized primarily as a
function of chlorophyll concentration \([Chl\,a]\).

\subsection{Phytoplankton Absorption}\label{phytoplankton-absorption}

The phytoplankton absorption coefficient \(a_{ph}\) follows the
formulation defined in the generalized model (see Section 6.3).

\subsection{Particulate Absorption}\label{particulate-absorption}

For open ocean waters, the contribution of non-algal particles (NAP) is
assumed negligible. Therefore the combined detrital and CDOM absorption
coefficient \(a_{dg}\) depends only on phytoplankton.

The value at 440 nm is parameterized as

\[
a_{dg}(440) = p_2\, a_p(440,[Chl\,a])
\]

where

\[
p_2 = 0.3 + \frac{5.7\,R_2\,a_p(440,[Chl\,a])}{0.02 + a_p(440,[Chl\,a])}.
\]

The parameter \(R_2\) is assumed to be

\[
R_2 = 0.5.
\]

The exponential spectral slope is fixed as

\[
S_{dg} = 0.018.
\]

\subsection{Particulate
Backscattering}\label{particulate-backscattering}

The particulate backscattering coefficient \(b_{bp}\) is also assumed to
depend only on phytoplankton in open waters.

The magnitude at 660 nm is parameterized as

\[
b_{bp}(660) = 0.347\,[Chl\,a]^{0.766}.
\]

The spectral slope of particulate backscattering is

\[
S_{bp} = -0.5\left(\log_{10}[Chl\,a] - 0.3\right),
\]

valid for

\[
0.02 < [Chl\,a] < 2 \; \text{mg m}^{-3}.
\]

Outside this range, the slope is set to zero.

\subsection{Backscattering Fraction}\label{backscattering-fraction}

The particulate backscattering fraction is parameterized as

\[
B_p = 0.002 + 0.01\left(0.5 - 0.25\log_{10}[Chl\,a]\right).
\]

The parameter \(B_p\) is assumed to be \textbf{spectrally independent}
(Huot et al. (2008)).

\subsection{Phase matrix (TBD)}\label{phase-matrix-tbd}

\begin{itemize}
\tightlist
\item
  water
\item
  NAP
\item
  Phytoplankton (FF)
\end{itemize}

\section{Bio-2 Optical Model (Five-Parameter
Model)}\label{bio-2-optical-model-five-parameter-model}

The \textbf{Bio-2 model} describes optically complex coastal waters
using a five-parameter parameterization.

The ocean is assumed to be homogeneous with a depth of \textbf{200 m}
and composed of four optical components:

\begin{enumerate}
\def\labelenumi{\arabic{enumi}.}
\tightlist
\item
  pure seawater
\item
  phytoplankton covariant particles
\item
  non-algal particles (NAP)
\item
  colored dissolved organic matter (CDOM)
\end{enumerate}

CDOM contributes only to absorption, while the other components
contribute to both absorption and scattering.

The absorption and scattering coefficients of pure seawater, \(a_w\) and
\(b_w\), are taken from laboratory measurements. The backscattering
fraction of pure water is assumed to be

\[
B_w = 0.5.
\]

\subsection{Phytoplankton Absorption}\label{phytoplankton-absorption-1}

Phytoplankton absorption is parameterized as

\[
a_{ph}(\lambda) = A_{ph}(\lambda)\,[Chl\,a]^{E_{ph}(\lambda)}.
\]

Here

\begin{itemize}
\tightlist
\item
  \(A_{ph}(\lambda)\) and \(E_{ph}(\lambda)\) are empirical coefficients
\item
  \([Chl\,a]\) is chlorophyll concentration in mg m\(^{-3}\).
\end{itemize}

\subsection{CDOM and Detrital
Absorption}\label{cdom-and-detrital-absorption}

The combined absorption coefficient of CDOM and NAP is modeled as

\[
a_{dg}(\lambda) =
a_{dg}(440)\exp\!\left[-S_{dg}(\lambda-440)\right].
\]

The parameter \(S_{dg}\) represents the exponential spectral slope.

\subsection{Particulate
Backscattering}\label{particulate-backscattering-1}

The particulate backscattering coefficient is modeled as

\[
b_{bp}(\lambda) =
b_{bp}(660)\left(\frac{\lambda}{660}\right)^{-S_{bp}}.
\]

The spectral slope \(S_{bp}\) controls the wavelength dependence.

\subsection{Backscattering Fraction}\label{backscattering-fraction-1}

The particulate backscattering fraction is

\[
B_p(\lambda) =
B_p(660)\left(\frac{\lambda}{660}\right)^{-S_{Bp}}.
\]

The parameter \(S_{Bp}\) defines the spectral variation of the
backscattering fraction.

\subsection{Total Inherent Optical
Properties}\label{total-inherent-optical-properties}

The total absorption coefficient is

\[
a(\lambda) =
a_w(\lambda) + a_{ph}(\lambda) + a_{dg}(\lambda).
\]

The total backscattering coefficient is

\[
b_b(\lambda) =
0.5\,b_w(\lambda) + b_{bp}(\lambda).
\]

These inherent optical properties are widely used to describe the ocean
color spectrum in bio-optical models.

\section{Bio-3 Optical Model (Three-Parameter
Model)}\label{bio-3-optical-model-three-parameter-model}

The \textbf{Bio-3 model} is a reduced-parameter version of the Bio-2
formulation designed to balance physical realism with retrieval
stability.

In this model, a subset of the Bio-2 parameters is retrieved while the
remaining parameters are constrained using empirical relationships
derived from observations.

\emph{(Details of the Bio-3 parameterization will be described soon.)}

\chapter{Multi-angle ocean color
product}\label{multi-angle-ocean-color-product}

Ocean color signals are often represented by the remote sensing
reflectance, \(R_{rs}(\lambda;\theta_0, \theta_v, \phi_v)\), which is
defined as the ratio of upwelling radiance originating from the ocean
body to downwelling irradiance just above the ocean surface, with units
of \(\mathrm{sr}^{-1}\) (Mobley et al. 2016). Here \(\theta_0\) and
\(\theta_v\) are the solar and viewing zenith angles, and \(\phi_v\) is
the relative viewing azimuth angle. For simplicity, the wavelength
dependence, \(\lambda\), which applies to all ocean reflectance-related
variables, is omitted in the following notation. To derive this
quantity, we need to remove the atmospheric path radiance and compute
the water-leaving reflectance (\(\rho_w\)) as the reflectance
originating from scattering in the ocean that reaches the satellite
sensor:

\begin{equation}\phantomsection\label{eq-rhow}{
\rho_w(\theta_0,\theta_v,\phi_v)
=
\rho_t(\theta_0,\theta_v,\phi_v)
-
\rho^{f}_{t,\mathrm{atm+sfc}}(\theta_0,\theta_v,\phi_v).
}\end{equation}

where \(\rho_t\) is the measured total reflectance at the top of
atmosphere (TOA), and \(\rho^{f}_{t,\mathrm{atm+sfc}}\) is the
reflectance with only atmospheric and ocean surface contributions (such
as sunglint, white caps, etc.) computed in the forward model, denoted by
superscript \(f\). With the multi-angle polarimetric observations
provided by PACE, wind speed can be retrieved simultaneously with
aerosol properties (M. Gao et al. 2021, 2022). The joint retrieval
algorithm enables the retrieval of ocean signals in regions affected by
sunglint, which are often excluded or difficult to process using
conventional single-angle ocean color remote sensing approaches (Harmel
and Chami 2013; Neukermans et al. 2018).

The remote sensing reflectance at arbitrary solar and viewing geometries
can be derived from the water-leaving reflectance by considering the
transmittance of both solar irradiance and surface radiance (M. Gao et
al. 2021) as:

\begin{equation}\phantomsection\label{eq-rrs}{
R_{rs}(\theta_0,\theta_v,\phi_v)
=
\left[
\frac{\rho_w(\theta_0,\theta_v,\phi_v)}{\pi}
\right]
\times
\left[
\frac{1}
{T_d(\theta_0)\,t_u(\theta_0,\theta_v,\phi_v)}
\right].
}\end{equation}

where \(T_d\) represents the downwelling irradiance transmittance of the
solar irradiance from the top of atmosphere to the surface, and \(t_u\)
represents the upwelling radiance transmittance for the water-leaving
radiance from the bottom of atmosphere to the satellite sensor, with
more details in M. Gao et al. (2021).

Furthermore, a BRDF correction procedure is commonly applied by mapping
\(R_{rs}(\theta_0, \theta_v, \phi_v)\) at arbitrary geometry to the sun
at zenith and viewing angle at nadir, denoted as \(R_{rs}(0,0)\). The
reduction of the solar and viewing geometry angular dependency allows
for geometrically consistent cross-instrument comparisons and validation
with in-situ data. Based on the radiative transfer model used in this
study, the BRDF correction factor can be computed through the modeled
\(R_{rs}^f\) as defined in Eqs. Equation~\ref{eq-rhow} and
Equation~\ref{eq-rrs} by replacing the real measurement by the forward
model simulation:

\begin{equation}\phantomsection\label{eq-cbrdf}{
C_{\mathrm{BRDF}}(\theta_0,\theta,\phi)
=
\frac{R^f_{rs}(0,0)}
{R^f_{rs}(\theta_0,\theta,\phi)}.
}\end{equation}

The angular distribution of the water-leaving light field, and thus this
correction term, depends on ocean parameters including Chl-a and wind
speed, as well as aerosol properties that modulate the diffuse skylight.
All of these parameters are retrieved simultaneously in the algorithm.
Note that this is an extension of the previous model in M. Gao et al.
(2021) for smaller viewing zenith angles, and has been applied in
synthetic PACE polarimeter data (M. Gao, Franz, et al. 2023; Aryal et
al. 2024; Zhang et al. 2025). Such a correction can be applied to
estimate \(R_{rs}(0,0)\):

\begin{equation}\phantomsection\label{eq-rrs00}{
R_{rs}(0,0)
=
R_{rs}(\theta_0,\theta,\phi)
\times
C_{\mathrm{BRDF}}(\theta_0,\theta,\phi).
}\end{equation}

To expedite this process, we further developed two neural networks which
compute \(1/(T_d t_u)\) and \(C_{\mathrm{BRDF}}/(T_d t_u)\) as
summarized in Appendix B, which can be used to efficiently produce
\(R_{rs}\) with and without BRDF correction.

Since HARP2 and SPEXone makes measurements at multiple viewing zenith
angles in the along-track direction), we can apply the atmospheric
correction to each of these and derive \(R_{rs}\) with a new angle
dimension. This satellite-based angular remote sensing reflectance
contains information on oceanic particles and potentially also
uncertainties from imperfect atmospheric correction. At first order, we
can compute the angular mean and standard deviation. The mean value can
be used to reduce random noise, while the standard deviation can be used
to evaluate angular variability.

Moreover, if the satellite-derived \(R_{rs}\) agrees with the modeled
\(R^f_{rs}\) from the forward model, assuming the bio-optical model
accurately represents the real water properties, we would expect the
angular variation of the \(R_{rs}\) to be reduced by the BRDF
correction. Therefore, BRDF correction applied on angular \(R_{rs}\) can
be used as an indicator of the closure of the radiative transfer model
and atmospheric correction.

Multiple BRDF correction schemes have been proposed in the ocean color
literature (Morel, Antoine, and Gentili 2002; Park and Ruddick 2005; Lee
et al. 2011; He et al. 2017; Twardowski and Tonizzo 2018; Zhang et al.
2025; D'Alimonte et al. 2025; Pitarch et al. 2025). The angular
\(R_{rs}\) observations provided by PACE from simultaneous multi-angle
measurements can be used to evaluate these BRDF correction schemes
across different water types, as well as angular closure analyses before
and after BRDF correction on a pixel-by-pixel basis. For more details
and validation results, please refer to Meng Gao et al. (2026).

\chapter{Land Surface Reflectance Model (Ross--Li BRDF and Polarized
BPDF)}\label{land-surface-reflectance-model-rossli-brdf-and-polarized-bpdf}

This document describes the mathematical formulation of the land surface
reflectance model implemented in the \textbf{PACE simulator}.
\textbf{Need proof read to be consistent with the code and formulation}

The model combines: - Ross--Li kernel-driven BRDF\\
- Polarized bidirectional reflectance distribution function (BPDF)\\
- Fresnel reflection physics

The output is a \textbf{4×4 Mueller matrix} describing polarized surface
reflection.

\section{Surface Reflectance Model}\label{surface-reflectance-model}

The bidirectional reflectance distribution matrix \(\mathbf{R}\) of land
surface converts the downwelling irradiance vector \(\mathbf{F}_s\) to
upwelling radiance vector \(\mathbf{I}\) by:

\[
\mathbf{I}(\theta_v,\phi_v)= \frac{1}{\pi}\mathbf{R}(\theta_s,\theta_v,\phi_r) \\, |\cos\theta_s| \\, \mathbf{F}_s(\theta_s,\phi_s) 
\]

where \(\theta\) and \(\phi\) are the zenith and azimuth angles,
respectively; the subscripts \(s\) and \(v\) denote the incident and
reflected directions, respectively; \(\phi_r=\phi_v-\phi_s\) is the
relative azimuth angle of the reflected and incident directions.

In the FastMAPOL algorithm the matrix \(\mathbf{R}\) is modeled by:

\[
\mathbf{R}(\theta_s,\theta_v,\phi_r) =
\left[
f_{iso}
+
f_{vol} K_{vol}(\theta_s,\theta_v,\phi_r)
+
f_{geo} K_{geo}(\theta_s,\theta_v,\phi_r)
\right]\mathbf{E}
+
B_{pol}\mathbf{K}_{pol}
\]

where \(\mathbf{E}\) is the 4 \(\times\) 4 null matrix except for
\(E_{11}=1\),

\[
\mathbf{E} =
\begin{bmatrix}
1 & 0 & 0 & 0 \\
0 & 0 & 0 & 0 \\
0 & 0 & 0 & 0 \\
0 & 0 & 0 & 0
\end{bmatrix}
\]

\(\mathbf{K}_{pol}\) is the polarization Kernel proportional to the
Fresnel reflection matrix.

\subsection{Scaled Parameter Formulation (Training
Representation)}\label{scaled-parameter-formulation-training-representation}

For parameter sampling (e.g., in neural network training), a scaled
formulation is used to reduce the number of free parameters and capture
main physics:

\[
\mathbf{R}(\theta_s,\theta_v,\phi_r) =
\left[
f_{iso}(\lambda)
\left(
1
+
k_{vol} K_{vol}
+
k_{geo} K_{geo}
\right)
\right]\mathbf{E}
+
B_{pol}\mathbf{K}_{pol}
\]

where - \(f_{iso}(\lambda)\) is spectrally dependent, and - \(k_{vol}\)
and \(k_{geo}\) are spectrally invariant

The relationship to the physical parameters is:

\[
f_{vol}(\lambda) = f_{iso}(\lambda) \times k_{vol}
\]

\[
f_{geo}(\lambda) = f_{iso}(\lambda) \times k_{geo}
\]

\subsection{Parameter Definitions}\label{parameter-definitions}

\begin{longtable}[]{@{}ll@{}}
\toprule\noalign{}
Parameter & Description \\
\midrule\noalign{}
\endhead
\bottomrule\noalign{}
\endlastfoot
\(\theta_s\) & solar zenith angle \\
\(\phi_s\) & solar azimuth angle \\
\(\theta_v\) & viewing zenith angle \\
\(\phi_v\) & viewing azimuth angle \\
\(\phi_r=\phi_v-\phi_s\) & relative azimuth angle \\
\(f_{iso}(\lambda)\) & isotropic reflectance coefficient \\
\(f_{vol}(\lambda)\) & volumetric scattering coefficient \\
\(f_{geo}(\lambda)\) & geometric scattering coefficient \\
\(k_{vol}\) & scaled volumetric scattering coefficient \\
\(k_{geo}\) & scaled geometric scattering coefficient \\
\(B_{pol}\) & polarization scaling parameter \\
\end{longtable}

\subsection{Geometry Definition}\label{geometry-definition}

Define cosine terms:

\[
\mu_s = \cos\theta_s
\]

\[
\mu_v = \cos\theta_v
\]

Sine terms:

\[
\sin\theta_s = \sqrt{1-\mu_s^2}
\]

\[
\sin\theta_v = \sqrt{1-\mu_v^2}
\]

Relative azimuth:

\[
\phi_r = \phi_v - \phi_s
\]

\subsection{RTSOS Geometry Convention}\label{rtsos-geometry-convention}

The model follows the \textbf{RTSOS light ray convention}, i.e., the
zenith and azimuth angles are associated with the directional vector of
the ligh ray.

\subsubsection{Geometry Definitions}\label{geometry-definitions}

\begin{longtable}[]{@{}ll@{}}
\toprule\noalign{}
Quantity & RTSOS definition \\
\midrule\noalign{}
\endhead
\bottomrule\noalign{}
\endlastfoot
solar zenith & \(\theta_s > 90^\circ\) \\
solar azimuth & \(\phi_s = 0\) \\
viewing zenith & \(\theta_v < 90^\circ\) \\
relative azimuth & \(\phi_r = \phi_v - \phi_s\) \\
\end{longtable}

\subsubsection{Scattering Planes}\label{scattering-planes}

\begin{longtable}[]{@{}ll@{}}
\toprule\noalign{}
Condition & Geometry \\
\midrule\noalign{}
\endhead
\bottomrule\noalign{}
\endlastfoot
\(\phi_r = 0\) & forward scattering \\
\(\phi_r = \pi\) & backscattering \\
\end{longtable}

\subsubsection{Scattering Angle}\label{scattering-angle}

The scattering angle between illumination and viewing directions is

\[
\cos\Theta =
\mu_s\mu_v +
\sqrt{1-\mu_s^2}\sqrt{1-\mu_v^2}\cos(\phi_r)
\]

The phase angle used in BRDF models is

\[
\Theta_{phase} = \pi - \Theta
\]

\section{Ross Volumetric Kernel}\label{ross-volumetric-kernel}

The volumetric kernel represents multiple scattering within vegetation
canopies.

\[
K_{vol} =
\frac{\left(\frac{\pi}{2}-\Theta_{phase}\right)\cos\Theta_{phase} + \sin\Theta_{phase}}
{|\mu_s| + |\mu_v|}-\frac{\pi}{4}
\]

Reference: Roujean et al.~(1992), JGR Atmosphere; Wanner et al.~(1995),
JGR.

\section{Li Geometric Kernel
(Li-Sparse-R)}\label{li-geometric-kernel-li-sparse-r}

Define tangent terms:

\[
|\tan\theta_s| = \frac{\sqrt{1-\mu_s^2}}{|\mu_s|}
\]

\[
|\tan\theta_v| = \frac{\sqrt{1-\mu_v^2}}{|\mu_v|}
\]

Projected angles:

\[
\theta'_s = \tan^{-1}(B/R \cdot |\tan\theta_s|)
\]

\[
\theta'_v = \tan^{-1}(B/R \cdot |\tan\theta_v|)
\]

\begin{longtable}[]{@{}ll@{}}
\toprule\noalign{}
Parameter & Meaning \\
\midrule\noalign{}
\endhead
\bottomrule\noalign{}
\endlastfoot
\(B/R\) & vertical-to-horizontal crown ratio \\
\end{longtable}

\subsection{Projected Phase Angle}\label{projected-phase-angle}

\[
\cos\xi =
\cos\theta'_s\cos\theta'_v +
\sin\theta'_s\sin\theta'_v\cos(\phi_r+\pi)
\]

\subsection{Projected Distance}\label{projected-distance}

\[
D =
\sqrt{
\tan^2\theta'_s +
\tan^2\theta'_v -
2\tan\theta'_s\tan\theta'_v\cos(\phi_r+\pi)
}
\]

Define:

\[
\sec\theta'_s = \frac{1}{\cos\theta'_s}
\]

\[
\sec\theta'_v = \frac{1}{\cos\theta'_v}
\]

\subsection{\texorpdfstring{Ancillary Function
}{Ancillary Function }}\label{ancillary-function}

\[
\cos t =
\frac{H}{B}\frac{
\sqrt{
D^2 + (\tan\theta'_s\tan\theta'_v\sin(\phi_r+\pi))^2
}}
{\sec\theta'_s + \sec\theta'_v}
\]

If

\[
|\cos t| > 1
\]

then

\[
O = 0
\]

Otherwise

\[
O =
\frac{1}{\pi}(t - \sin t \\, \cos t)
(\sec\theta'_s + \sec\theta'_v)
\]

\section{Final Geometric Kernel}\label{final-geometric-kernel}

\[
K_{geo} = O - \sec\theta'_s - \sec\theta'_v + \frac{1}{2}(1+\cos\xi) \\, \sec\theta'_s \\,  \sec\theta'_v
\]

\section{Polarized Surface Reflectance
(BPDF)}\label{polarized-surface-reflectance-bpdf}

Scattering angle have been defined above \(\Theta\). Assuming a flat
surface which produce the same scattering angle, the incident angle
relative to the flat surface would be:

\[
\theta_1 =\frac{\pi - \Theta}{2}
\]

which also equivalent to \(\theta_1=(\theta_s+\theta_v)/2\)

\subsection{Fresnel Reflection}\label{fresnel-reflection}

Let the refractive index be

\[
n = n_r + i n_i
\]

Parallel polarization:

\[
r_p =
\frac{n^2\cos\theta_1 - \sqrt{n^2-\sin^2\theta_1}}
{n^2\cos\theta_1 + \sqrt{n^2-\sin^2\theta_1}}
\]

Perpendicular polarization:

\[
r_s =
\frac{\cos\theta_1 - \sqrt{n^2-\sin^2\theta_1}}
{\cos\theta_1 + \sqrt{n^2-\sin^2\theta_1}}
\]

\textbf{Assumption: \(n = 1.5 + 0i\) }

\subsection{Fresnel Mueller Matrix}\label{fresnel-mueller-matrix}

\[
A = \frac{1}{2}|r_p|^2
\]

\[
B = \frac{1}{2}|r_s|^2
\]

\[
\gamma = r_p^* r_s
\]

\[
\mathbf{F} =
\begin{bmatrix}
A+B & A-B & 0 & 0 \\
A-B & A+B & 0 & 0 \\
0 & 0 & Re(\gamma) & -Im(\gamma) \\
0 & 0 & Im(\gamma) & Re(\gamma)
\end{bmatrix}
\]

Nadir reflectance (A+B) can be approximated as \([(n-1)/(n+1)]^2=0.04\)
for n=1.5. Brewster angle can be observed at the maximum of
Abs(A-B)/(A+B).

\subsection{Polarization Kernel}\label{polarization-kernel}

In the current framework, we adopt the polarized Kernel based on Maignan
et al., (2009):

\[
\mathbf{K}_{pol} =
\frac{
e^{-\tan\theta_1}e^{-v_{fac}}
}{4(|\mu_s|+|\mu_v|)}
\mathbf{L(\pi-i_2)}\mathbf{F} (n, \theta_1) \mathbf{L(i_1)}
\]

where \(\mathbf{L(i_1)}\) and \(\mathbf{L(\pi-i_2)}\) rotate the Stokes
vector from the incident meridian plane to the scattering plane and back
to the viewing meridian plane; \(i_1\) and \(i_2\) are the rotation
angles; The factor \(v_{fac}\) is:

\[
v_{fac} = 0.1
\]

For scalar radiative transfer simulations, only \(R_{11}\) is used.

\begin{center}\rule{0.5\linewidth}{0.5pt}\end{center}

\section{References}\label{references}

Maignan, F., et al., (2009) \emph{Polarized reflectances of natural
surfaces: Spaceborne measurements and analytical modeling.} Remote
Sensing of Environment, Volume 113, Issue 12, Pages 2642-2650,
https://doi.org/10.1016/j.rse.2009.07.022.

MODIS Surface Reflectance ATBD\\
https://modis.gsfc.nasa.gov/data/atbd/atbd\_mod09.pdf

Roujean, J.-L., M.Leroy, and P.-Y.Deschamps (1992), A bidirectional
reflectance model of the Earth's surface for the correction of remote
sensing data, J. Geophys. Res., 97(D18), 20455--20468,
doi:10.1029/92JD01411.

Wanner, W., Li, X., \& Strahler, A. (1995)\\
\emph{On the derivation of kernels for kernel-driven BRDF models.}\\
Journal of Geophysical Research.

\phantomsection\label{refs}
\begin{CSLReferences}{1}{0}
\bibitem[\citeproctext]{ref-Anderson:1986aa}
Anderson, G., Shepard Clough, F. Kneizys, J. Chetwynd, and Eric Shettle.
1986. \emph{AFGL Atmospheric Constituent Profiles (0.120km)}.
AFGL-TR-86-0110. Air Force Geophysics Lab., Hanscom AFB, MA (USA).

\bibitem[\citeproctext]{ref-Aryal:2026aa}
Aryal, Kamal, Meng Gao, Peng-Wang Zhai, Bryan A. Franz, Kirk
Knobelspiesse, Amir Ibrahim, and Jeremy Werdell. 2026. {``Aerosol and
Ocean Color Retrievals from FastMAPOL/Component Algorithm Using NASA's
PACE Polarimetric Measurements.''} \emph{Opt. Express} 34 (16):
29139--55. \url{https://doi.org/10.1364/OE.607432}.

\bibitem[\citeproctext]{ref-Aryal:2024aa}
Aryal, Kamal, Peng-Wang Zhai, Meng Gao, Bryan A. Franz, Kirk
Knobelspiesse, and Yongxiang Hu. 2024. {``Machine Learning Based Aerosol
and Ocean Color Joint Retrieval Algorithm for Multiangle Polarimeters
over Coastal Waters.''} \emph{Opt. Express} 32 (17): 29921--42.
\url{https://doi.org/10.1364/OE.522794}.

\bibitem[\citeproctext]{ref-Baydin:2018aa}
Baydin, Atilim Gunes, Barak A. Pearlmutter, Alexey Andreyevich Radul,
and Jeffrey Mark Siskind. 2018. {``Automatic Differentiation in Machine
Learning: A Survey.''} \emph{Journal of Machine Learning Research} 18:
1--43. \url{https://arxiv.org/abs/1502.05767}.

\bibitem[\citeproctext]{ref-Branch:1999aa}
Branch, Mary Ann, Thomas F. Coleman, and Yuying Li. 1999. {``A Subspace,
Interior, and Conjugate Gradient Method for Large-Scale
Bound-Constrained Minimization Problems.''} \emph{SIAM Journal on
Scientific Computing} 21 (1): 1--23.
\url{https://doi.org/10.1137/S1064827595289108}.

\bibitem[\citeproctext]{ref-DAlimonte:2025aa}
D'Alimonte, Davide, Tamito Kajiyama, Jaime Pitarch, Vittorio Ernesto
Brando, Marco Talone, Constant Mazeran, Michael Twardowski, et al. 2025.
{``Comparison of Correction Methods for Bidirectional Effects in Ocean
Colour Remote Sensing.''} \emph{Remote Sensing of Environment} 321:
114606. https://doi.org/\url{https://doi.org/10.1016/j.rse.2025.114606}.

\bibitem[\citeproctext]{ref-Duan:2005aa}
Duan, Minzheng, Qilong Min, and Jiangnan Li. 2005. {``A Fast Radiative
Transfer Model for Simulating High-Resolution Absorption Bands.''}
\emph{Journal of Geophysical Research: Atmospheres} 110 (D15201).
https://doi.org/\url{https://doi.org/10.1029/2004JD005590}.

\bibitem[\citeproctext]{ref-Dubovik:2011aa}
Dubovik, O., M. Herman, A. Holdak, T. Lapyonok, D. Tanré, J. L. Deuzé,
F. Ducos, A. Sinyuk, and A. Lopatin. 2011. {``Statistically Optimized
Inversion Algorithm for Enhanced Retrieval of Aerosol Properties from
Spectral Multi-Angle Polarimetric Satellite Observations.''}
\emph{Atmospheric Measurement Techniques} 4 (5): 975--1018.
\url{https://doi.org/10.5194/amt-4-975-2011}.

\bibitem[\citeproctext]{ref-Dubovik:2006aa}
Dubovik, Oleg, Alexander Sinyuk, Tatyana Lapyonok, Brent N. Holben,
Michael Mishchenko, Ping Yang, Tom F. Eck, et al. 2006. {``Application
of Spheroid Models to Account for Aerosol Particle Nonsphericity in
Remote Sensing of Desert Dust.''} \emph{Journal of Geophysical Research:
Atmospheres} 111: D11208. \url{https://doi.org/10.1029/2005JD006619}.

\bibitem[\citeproctext]{ref-Fournier:1994aa}
Fournier, Georges R., and J. Luc Forand. 1994. {``Analytic Phase
Function for Ocean Water.''} \emph{Proc.SPIE} 2258: 2258 - 2258 - 8.
\url{https://doi.org/10.1117/12.190063}.

\bibitem[\citeproctext]{ref-Gao:2026aa}
Gao, Meng, Kamal Aryal, Peng-Wang Zhai, Kirk Knobelspiesse, Bryan Franz,
Brian Cairns, Ivona Cetinić, et al. 2026. {``Where Sky Meets Sea:
Integrated Aerosol and Ocean Color Retrieval from PACE SPEXone
Multi-Angle Polarimetry.''} \emph{Remote Sensing of Environment} 345:
115603. https://doi.org/\url{https://doi.org/10.1016/j.rse.2026.115603}.

\bibitem[\citeproctext]{ref-Gao:2021bb}
Gao, Meng, Kirk Knobelspiesse, Bryan A. Franz, Peng-Wang Zhai, Vanderlei
Martins, Sharon P. Burton, Brian Cairns, et al. 2021. {``Adaptive Data
Screening for Multi-Angle Polarimetric Aerosol and Ocean Color Remote
Sensing Accelerated by Deep Learning.''} \emph{Frontiers in Remote
Sensing} 2: 46. \url{https://doi.org/10.3389/frsen.2021.757832}.

\bibitem[\citeproctext]{ref-Gao:2018aa}
Gao, Meng, Peng-Wang Zhai, Bryan Franz, Yongxiang Hu, Kirk
Knobelspiesse, P. Jeremy Werdell, Amir Ibrahim, Feng Xu, and Brian
Cairns. 2018. {``Retrieval of Aerosol Properties and Water-Leaving
Reflectance from Multi-Angular Polarimetric Measurements over Coastal
Waters.''} \emph{Opt. Express} 26 (7): 8968--89.
\url{https://doi.org/10.1364/OE.26.008968}.

\bibitem[\citeproctext]{ref-Gao:2021aa}
Gao, M., B. A. Franz, K. Knobelspiesse, P.-W. Zhai, V. Martins, S.
Burton, B. Cairns, et al. 2021. {``Efficient Multi-Angle Polarimetric
Inversion of Aerosols and Ocean Color Powered by a Deep Neural Network
Forward Model.''} \emph{Atmospheric Measurement Techniques} 14 (6):
4083--4110. \url{https://doi.org/10.5194/amt-14-4083-2021}.

\bibitem[\citeproctext]{ref-Gao:2023aa}
Gao, M., B. A. Franz, P.-W. Zhai, K. Knobelspiesse, A. M. Sayer, X. Xu,
J. V. Martins, et al. 2023. {``Simultaneous Retrieval of Aerosol and
Ocean Properties from PACE HARP2 with Uncertainty Assessment Using
Cascading Neural Network Radiative Transfer Models.''} \emph{Atmospheric
Measurement Techniques} 16 (23): 5863--81.
\url{https://doi.org/10.5194/amt-16-5863-2023}.

\bibitem[\citeproctext]{ref-Gao:2023bb}
Gao, M., K. Knobelspiesse, B. A. Franz, P.-W. Zhai, B. Cairns, X. Xu,
and J. V. Martins. 2023. {``The Impact and Estimation of Uncertainty
Correlation for Multi-Angle Polarimetric Remote Sensing of Aerosols and
Ocean Color.''} \emph{Atmospheric Measurement Techniques} 16 (8):
2067--87. \url{https://doi.org/10.5194/amt-16-2067-2023}.

\bibitem[\citeproctext]{ref-Gao:2022aa}
Gao, M., K. Knobelspiesse, B. A. Franz, P.-W. Zhai, A. M. Sayer, A.
Ibrahim, B. Cairns, et al. 2022. {``Effective Uncertainty Quantification
for Multi-Angle Polarimetric Aerosol Remote Sensing over Ocean.''}
\emph{Atmospheric Measurement Techniques} 15 (16): 4859--79.
\url{https://doi.org/10.5194/amt-15-4859-2022}.

\bibitem[\citeproctext]{ref-Gelaro:2017aa}
Gelaro, Ronald, Will McCarty, Max J. Suárez, Ricardo Todling, Andrea
Molod, Lawrence Takacs, Cynthia A. Randles, et al. 2017. {``The
Modern-Era Retrospective Analysis for Research and Applications, Version
2 (MERRA-2).''} \emph{Journal of Climate} 30 (14): 5419--54.
\url{https://doi.org/10.1175/JCLI-D-16-0758.1}.

\bibitem[\citeproctext]{ref-Harmel:2013aa}
Harmel, Tristan, and Malik Chami. 2013. {``Estimation of the Sunglint
Radiance Field from Optical Satellite Imagery over Open Ocean:
Multidirectional Approach and Polarization Aspects.''} \emph{Journal of
Geophysical Research: Oceans} 118 (1): 76--90.
https://doi.org/\url{https://doi.org/10.1029/2012JC008221}.

\bibitem[\citeproctext]{ref-He:2017aa}
He, Shuangyan, Xiaodong Zhang, Yuanheng Xiong, and Deric Gray. 2017.
{``A Bidirectional Subsurface Remote Sensing Reflectance Model
Explicitly Accounting for Particle Backscattering Shapes.''}
\emph{Journal of Geophysical Research: Oceans} 122 (11): 8614--26.
https://doi.org/\url{https://doi.org/10.1002/2017JC013313}.

\bibitem[\citeproctext]{ref-Huot:2008aa}
Huot, Y, A Morel, M S Twardowski, D Stramski, and R A Reynolds. 2008.
{``{Particle optical backscattering along a chlorophyll gradient in the
upper layer of the eastern South Pacific Ocean}.''}
\emph{Biogeosciences} 5 (2): 495--507.

\bibitem[\citeproctext]{ref-Frederick:2006aa}
James, Frederick. 2006. \emph{Statistical Methods in Experimental
Physics}. World Scientific, Singapore.

\bibitem[\citeproctext]{ref-Knobelspiesse:2012aa}
Knobelspiesse, Kirk, Brian Cairns, Michael Mishchenko, Jacek Chowdhary,
Kostas Tsigaridis, Bastiaan van Diedenhoven, William Martin, Matteo
Ottaviani, and Mikhail Alexandrov. 2012. {``Analysis of Fine-Mode
Aerosol Retrieval Capabilities by Different Passive Remote Sensing
Instrument Designs.''} \emph{Optics Express} 20 (19): 21457--84.
\url{https://doi.org/10.1364/OE.20.021457}.

\bibitem[\citeproctext]{ref-Knobelspiesse:2019aa}
Knobelspiesse, Kirk, Qian Tan, Carol Bruegge, Brian Cairns, Jacek
Chowdhary, Bastiaan van Diedenhoven, David Diner, et al. 2019.
{``Intercomparison of Airborne Multi-Angle Polarimeter Observations from
the Polarimeter Definition Experiment.''} \emph{Appl. Opt.} 58 (3):
650--69. \url{https://doi.org/10.1364/AO.58.000650}.

\bibitem[\citeproctext]{ref-Koepke:1984aa}
Koepke, P. 1984. {``Effective Reflectance of Oceanic Whitecaps.''}
\emph{Appl. Opt.} 23: 1816--24.
\href{https://doi.org/doi:\%2010.1364/AO.23.001816}{https://doi.org/doi:
10.1364/AO.23.001816}.

\bibitem[\citeproctext]{ref-Kokhanovsky:2003aa}
Kokhanovsky, Alexander A. 2003. {``{Parameterization of the Mueller
matrix of oceanic waters}.''} \emph{Journal of Geophysical Research:
Oceans} 108 (C6): 3175.

\bibitem[\citeproctext]{ref-Lee:2011aa}
Lee, Zhong Ping, Keping Du, Kenneth J. Voss, Giuseppe Zibordi, Bertrand
Lubac, Robert Arnone, and Alan Weidemann. 2011. {``An
Inherent-Optical-Property-Centered Approach to Correct the Angular
Effects in Water-Leaving Radiance.''} \emph{Appl. Opt.} 50 (19):
3155--67. \url{https://doi.org/10.1364/AO.50.003155}.

\bibitem[\citeproctext]{ref-Mobley:2016aa}
Mobley, Curtis D., Jeremy Werdell, Bryan Franz, Ziauddin Ahmad, and Sean
Bailey. 2016. \emph{Atmospheric Correction for Satellite Ocean Color
Radiometry}. National Aeronautics; Space Administration, Washington,
D.C.United States.
\url{https://oceancolor.gsfc.nasa.gov/docs/technical/NASA-TM-2016-217551.pdf}.

\bibitem[\citeproctext]{ref-Morel:2002aa}
Morel, André, David Antoine, and Bernard Gentili. 2002. {``Bidirectional
Reflectance of Oceanic Waters: Accounting for Raman Emission and Varying
Particle Scatter~Ing Phase Function.''} \emph{Appl. Opt.} 41 (30):
6289--6306. https://doi.org/\url{https://doi.org/10.1364/AO.41.006289}.

\bibitem[\citeproctext]{ref-Neukermans:2018aa}
Neukermans, Griet, Tristan Harmel, Martí Galí, Natalia Rudorff, Jacek
Chowdhary, Oleg Dubovik, Chris Hostetler, et al. 2018. {``Harnessing
Remote Sensing to Address Critical Science Questions on Ocean-Atmosphere
Interactions.''} \emph{Elementa: Science of the Anthropocene} 6
(November): 71. \url{https://doi.org/10.1525/elementa.331}.

\bibitem[\citeproctext]{ref-Park:2005aa}
Park, Young-Je, and Kevin Ruddick. 2005. {``Model of Remote-Sensing
Reflectance Including Bidirectional Effects for Case 1 and Case 2
Waters.''} \emph{Appl. Opt.} 44 (7): 1236--49.
\url{https://doi.org/10.1364/AO.44.001236}.

\bibitem[\citeproctext]{ref-Paszke:2019aa}
Paszke, Adam, Sam Gross, Francisco Massa, Adam Lerer, James Bradbury,
Gregory Chanan, Trevor Killeen, et al. 2019. {``PyTorch: An Imperative
Style, High-Performance Deep Learning Library.''} Edited by H. Wallach,
H. Larochelle, A. Beygelzimer, F. dAlché-Buc, E. Fox, and R. Garnett.
\emph{Advances in Neural Information Processing Systems 32}, 8024--35.
\url{http://papers.neurips.cc/paper/9015-pytorch-an-imperative-style-high-performance-deep-learning-library.pdf}.

\bibitem[\citeproctext]{ref-Pitarch:2025aa}
Pitarch, Jaime, Vittorio Ernesto Brando, Marco Talone, Constant Mazeran,
Davide D'Alimonte, Tamito Kajiyama, Ewa Kwiatkowska, David Dessailly,
and Juan Ignacio Gossn. 2025. {``Analytical Modeling and Correction of
the Ocean Colour Bidirectional Reflectance Across Water Types.''}
\emph{Remote Sensing of Environment}.
https://doi.org/\url{https://doi.org/10.1016/j.rse.2025.114920}.

\bibitem[\citeproctext]{ref-Priestley:1983aa}
Priestley, M. B. 1983. \emph{Spectral Analysis and Time Series}.
Academic Press.

\bibitem[\citeproctext]{ref-Rodgers:2000aa}
Rodgers, C. D. 2000. \emph{Inverse Methods for Atmospheric
Sounding:theory and Practice}. World Scientific World Scientific
Publishing, Singapore.
https://doi.org/\url{https://doi.org/10.1142/3171}.

\bibitem[\citeproctext]{ref-Sayer:2020aa}
Sayer, A. M., Y. Govaerts, P. Kolmonen, A. Lipponen, M. Luffarelli, T.
Mielonen, F. Patadia, et al. 2020. {``A Review and Framework for the
Evaluation of Pixel-Level Uncertainty Estimates in Satellite Aerosol
Remote Sensing.''} \emph{Atmospheric Measurement Techniques} 13 (2):
373--404. \url{https://doi.org/10.5194/amt-13-373-2020}.

\bibitem[\citeproctext]{ref-Seabold:2010aa}
Seabold, Skipper, and Josef Perktold. 2010. {``Statsmodels: Econometric
and Statistical Modeling with Python.''} In \emph{Proceedings of the 9th
Python in Science Conference}.

\bibitem[\citeproctext]{ref-Twardowski:2018aa}
Twardowski, Michael, and Alberto Tonizzo. 2018. {``Ocean Color
Analytical Model Explicitly Dependent on the Volume Scattering
Function.''} \emph{Applied Sciences} 8 (12).
\url{https://doi.org/10.3390/app8122684}.

\bibitem[\citeproctext]{ref-Virtanen:2018aa}
Virtanen, T. H., P. Kolmonen, L. Sogacheva, E. Rodrı́guez, G. Saponaro,
and G. de Leeuw. 2018. {``Collocation Mismatch Uncertainties in
Satellite Aerosol Retrieval Validation.''} \emph{Atmospheric Measurement
Techniques} 11 (2): 925--38.
\url{https://doi.org/10.5194/amt-11-925-2018}.

\bibitem[\citeproctext]{ref-Voss:1984aa}
Voss, Kenneth J, and Edward S Fry. 1984. {``{Measurement of the Mueller
matrix for ocean water}.''} \emph{Applied Optics} 23 (23): 4427--39.

\bibitem[\citeproctext]{ref-Xu:2016aa}
Xu, F, O Dubovik, P W Zhai, D J Diner, O V Kalashnikova, F C Seidel, P
Litvinov, et al. 2016. {``{Joint retrieval of aerosol and water-leaving
radiance from multispectral, multiangular and polarimetric measurements
over ocean}.''} \emph{Atmospheric Measurement Techniques} 9 (7):
2877--907.

\bibitem[\citeproctext]{ref-Zhai:2022aa}
Zhai, Peng-Wang, and Yongxiang Hu. 2022. {``An Improved Pseudo Spherical
Shell Algorithm for Vector Radiative Transfer.''} \emph{Journal of
Quantitative Spectroscopy and Radiative Transfer} 282: 108132.
https://doi.org/\url{https://doi.org/10.1016/j.jqsrt.2022.108132}.

\bibitem[\citeproctext]{ref-Zhai:2010aa}
Zhai, Peng-Wang, Yongxiang Hu, Jacek Chowdhary, Charles R Trepte,
Patricia L Lucker, and Damien B Josset. 2010. {``{A vector radiative
transfer model for coupled atmosphere and ocean systems with a rough
interface}.''} \emph{Journal of Quantitative Spectroscopy and Radiative
Transfer} 111 (7): 1025--40.

\bibitem[\citeproctext]{ref-Zhai:2009aa}
Zhai, Peng-Wang, Yongxiang Hu, Charles R Trepte, and Patricia L Lucker.
2009. {``{A vector radiative transfer model for coupled atmosphere and
ocean systems based on successive order of scattering method}.''}
\emph{Optics Express} 17 (4): 2057--79.

\bibitem[\citeproctext]{ref-Zhai:2015aa}
Zhai, Peng-Wang, Yongxiang Hu, David M Winker, Bryan A Franz, and
Emmanuel Boss. 2015. {``{Contribution of Raman scattering to polarized
radiation field in ocean waters}.''} \emph{Optics Express} 23 (18):
23582--96.

\bibitem[\citeproctext]{ref-Zhai:2017aa}
Zhai, Peng-Wang, Kirk Knobelspiesse, Amir Ibrahim, Bryan A Franz,
Yongxiang Hu, Meng Gao, and Robert Frouin. 2017. {``{Water-leaving
contribution to polarized radiation field over ocean}.''} \emph{Optics
Express} 25 (16): A689--708.
https://doi.org/\url{DOI:https://doi.org/10.1364/OE.25.00A689}.

\bibitem[\citeproctext]{ref-Zhang:2025aa}
Zhang, Xiaodong, Meng Gao, Shuangyan He, and Lucas Barbedo. 2025.
{``Deriving Anisotropic Correction for Upwelling Radiance from PACE's
Multi-Angle Polarimetry.''} \emph{Remote Sensing of Environment} 320:
114647. https://doi.org/\url{https://doi.org/10.1016/j.rse.2025.114647}.

\end{CSLReferences}




\end{document}
